\chapimg{fig/sh14.jpg}
\IfFileExists{extra/lead14.tex}{\chapterlead{\subhd{After this chapter you can} explain why drainage and salinity must be treated together; compute Hooghoudt's drain spacing with the equivalent depth; compute the leaching requirement and the leaching fraction; and design a subsurface drainage system with drain length and pipe size.}
}{}
\chapter{Subsurface Drainage Design}\label{chap:14}
\markboth{Subsurface Drainage Design}{Subsurface Drainage Design}
\IfFileExists{extra/bg14.tex}{\section{Introduction}
Irrigation adds water and salts. Where the water table is shallow, or the soil drains slowly, the salts stay in the root zone and the water table rises, so the crop is both too wet and too salty. Drainage removes the surplus water and, with it, the salt; it is the part of the water balance that keeps irrigation sustainable. This chapter gives the two calculations of drainage design: the spacing of subsurface drains and the leaching needed to keep the salinity at the tolerable level.

\section{Types of drainage}
Drainage is \emph{surface} (ditches, land levelling, bedding) or \emph{subsurface} (buried perforated pipes or tiles, mole drains, tube wells for vertical drainage). Subsurface pipe drains suit flat irrigated land with fine texture and a barrier below.

\section{Steady-state spacing: the Hooghoudt equation}
For steady recharge $q$ onto a water table midway between two drains of spacing $L$, with the water table $h$ above drain level, the Hooghoudt equation is

\begin{equation}q=\frac{8K_2d_e\,h+4K_1h^2}{L^2},\end{equation}


where $K_1$ is the conductivity above the drains, $K_2$ below, and $d_e$ the equivalent depth that allows for the convergence of flow to the drains: $d_e=d/\bigl[1+\frac{d}{L}\bigl(\frac{8}{\pi}\ln\frac{d}{r_0}-3.4\bigr)\bigr]$ for $d/L<0.3$, where $d$ is the depth of the barrier below the drains and $r_0$ the drain radius. Since $d_e$ depends on $L$, the equation is solved by iteration.

\section{Salinity and the leaching requirement}
A crop tolerates a saturation-extract salinity $EC_e$. The water that drains below the roots carries away the salt of the irrigation water, concentrated by evapotranspiration. The \kw{leaching requirement} is $LR=EC_w/(5EC_e-EC_w)$. For irrigation water of $EC_w=1.5$ dS\,m$^{-1}$ and a crop with $EC_e=3.0$: $LR=0.111$, so 11.1 per cent of the applied water must drain.

}{}
\IfFileExists{extra/th14.tex}{\section{The need for drainage}
A soil is waterlogged when the water table is so close to the surface that the pores in the root zone are filled with water and the roots cannot breathe. The causes are heavy rain, irrigation with little drainage, seepage from canals and a low permeability of the subsoil. The effects are the lack of oxygen (the roots die within a few days in warm soil), the loss of nitrogen by denitrification, a cold seedbed, and, in dry regions, the rise of salt by capillary action from the water table. Drainage removes the excess water, lowers the water table and allows the leaching of salts.

\section{Kinds of drainage}
\textbf{Surface drainage} removes the water that ponds on the surface, by shaping the land and by open ditches. \textbf{Subsurface drainage} lowers the water table by pipes buried below the surface or by open drains deep enough to reach it. The layouts of subsurface drains are:
\begin{exlist}
\item \kw{Random} (natural) system: pipes in the wet spots only.
\item \kw{Parallel} system: laterals at equal spacing, which feed a collector; for flat, uniform land.
\item \kw{Herringbone}: laterals join the collector at an angle, from both sides; for narrow depressions.
\item \kw{Interceptor} (cut-off) drain: a line across a slope to catch seepage from higher land.
\end{exlist}
The pipes are of corrugated plastic or clay, 60 to 100 mm in diameter for laterals, with slots; they are wrapped in an \kw{envelope} (gravel, geotextile) to prevent fine soil from entering. Spacings of 10 to 30 m and depths of 1.0 to 1.5 m are common.

\section{Steady-state and non-steady-state design}
In the \kw{steady-state} method the recharge (rain or irrigation) is taken as constant, and the water table is assumed to be at a fixed height at midway between drains; the equation of Hooghoudt gives the spacing. In the \kw{non-steady-state} method (Glover and Dumm), the recharge is intermittent and the water table falls after each event, so that the spacing is found from the time in which the water table must fall by a given amount. Non-steady design is more realistic in rain-fed humid areas, steady design in irrigated areas with a constant recharge; the two agree when the recharge is frequent.

\section{Assumptions of the Hooghoudt equation}
The equation is derived with the assumptions of Dupuit and Forchheimer: (1) the flow to the drain is horizontal, and the hydraulic gradient equals the slope of the water table; (2) the soil is homogeneous and isotropic; (3) the drains reach the impermeable layer, or an \emph{equivalent depth} $d$ replaces the actual depth to account for the convergence of the flow near the drain. The equivalent depth is less than the actual depth to the barrier; and when the barrier is deep, the flow below the drain is large, and the radial resistance is the main part of the loss.

\section{Drainage coefficient and the design discharge}
The \kw{drainage coefficient} $q_d$ is the depth of water to be removed per day, in mm\,d$^{-1}$. It is taken from the climate and the crop: 1 to 3 mm\,d$^{-1}$ for irrigated land in arid regions, 5 to 10 for humid regions, and up to 20 for vegetables. A larger coefficient needs closer drains and is more costly; the aim is to keep the water table below the depth that the crop can tolerate for the time it can tolerate it.

\section{Drainage and the environment}
Drainage removes water, but also nitrate, phosphorus and pesticides, and the drain water may harm the water bodies that receive it. Practices to reduce the loss are \emph{controlled drainage} (a weir in the outlet keeps the water table higher between rains), \emph{shallower drains}, the use of buffer strips and wetlands at the outlet, and a better timing of fertiliser. The nitrate concentration of the drain water (Chapter~\ref{chap:17}) is compared with the limit of 11.3 mg N per litre.
}{}
\IfFileExists{extra/deep14.tex}{\section{Drainage: what the layer does}
The drainage layer answers three questions of the farm that is irrigated: how far apart must the drains be to keep the water table below the roots; how large must the pipes be to carry the water away; and how much extra water must be applied, and drained, to wash out the salt. It takes the recharge of the water table from the field water balance (the deep percolation of rain and the loss of the irrigation design), the conductivity of the soil from the Saxton and Rawls functions, and the crop's tolerance of salinity and of a shallow water table.

\section{Types of drains}
The application offers three. A \textbf{buried pipe drain} is a perforated corrugated pipe in a gravel envelope; the flow enters through the perforations, the pipe runs half full, and the effective wetted perimeter is $u=\pi r_0$ with $r_0=0.05$ m (a 100-mm pipe). A \textbf{pipe drain with a control gate} has an adjustable outlet that holds the water table at a chosen height $g$ above the drain level, so that in the dry months the water is kept in the soil for the crop; it is the \emph{controlled drainage} of the application. An \textbf{open ditch} is an excavated channel with side slope 1.5 and bottom width $b_w$; the wetted perimeter is $u=b_w+2D_w\sqrt{1+1.5^2}$ with a water depth $D_w=0.3$ m, Manning $n=0.03$, and it takes more land (the ditch occupies $b_w+1.5\,\text{depth}$ of the width) and is cleaned each year. Pipe drains use $n=0.016$ (corrugated polyethylene).

\begin{table}[!htbp]\centering\footnotesize\setlength{\tabcolsep}{3pt}\caption{The drain types of the application.}
\begin{tabular}{llll}\toprule
\hd{} & \hd{Buried pipe} & \hd{Pipe with gate} & \hd{Open ditch}\\\midrule
Wetted perimeter $u$ & $\pi r_0=0.157$ m & $\pi r_0=0.157$ m & $b_w+2D_w\sqrt{1+z^2}$\\
Manning $n$ & 0.016 & 0.016 & 0.030\\
Default depth of drain (m) & 1.2 & 1.2 & 1.2\\
Controls water table & no & yes (gate $g$) & no\\
Land lost & none & none & ditch and spoil bank\\
Maintenance & low & low & high (annual cleaning)\\\bottomrule\end{tabular}\end{table}

\section{Steady design: Hooghoudt and the equivalent depth}
When the recharge $q$ is steady, the water table between two parallel drains of spacing $L$ rises to a height $h$ above the drain level in the middle; Hooghoudt's equation relates them:

\begin{equation}\boxed{q=\frac{8K_2\,d_e\,h+4K_1\,h^{2}}{L^{2}},\qquad L=\sqrt{\frac{8K_2d_eh+4K_1h^2}{q}},}\label{eq:dr-h}\end{equation}
\par\noindent where $q$ is the design discharge (recharge) (m\,d$^{-1}$), $L$ the spacing of the drains (m), $h$ the allowed height of the water table midway above the drain level (m), $K_1$ the saturated conductivity of the layer above the drain level and $K_2$ that of the layer below it (m\,d$^{-1}$) and $d_e$ the equivalent depth (m). In the application $K_1=K_s$ of the soil (Saxton and Rawls) and $K_2=K_1\times K_{sub}$ with $K_{sub}$ chosen by the user (default 1).



Close to the drain the flow lines converge, which lengthens the path and raises the head needed; Hooghoudt replaced the real depth $D$ of the barrier below the drains by a smaller \emph{equivalent depth} $d_e$ that allows for this. The application computes it from the exact solution of Moody (van Beers' form):

\begin{equation}\boxed{d_e=\min\left\{D,\ \frac{\pi L/8}{\ln(L/u)+F(x)}\right\},\quad F(x)=\sum_{n=1,3,5,\ldots,29}\frac{4e^{-2nx}}{n\,(1-e^{-2nx})},\quad x=\frac{2\pi D}{L},}\label{eq:dr-de}\end{equation}
\par\noindent where $d_e$ is the equivalent depth (m), $D$ the depth of the impermeable layer below the drains (m), $L$ the spacing (m), $u$ the wetted perimeter of the drain (m), $x$ a dimensionless number and $F(x)$ the series (15 terms) that tends to zero when the barrier is deep ($x$ large) and corrects for the finite depth otherwise.



Since $d_e$ depends on $L$ the equation is solved by iteration: start with $L=\sqrt{(4K_1h^2+8K_2Dh)/q}$, compute $d_e(L)$, update $L$, and average the old and the new value until the change is below 1 cm. The critical depth of the water table is $z_c=\mathrm{clip}(0.3+0.4R_{max}/100,\ 0.4,\ 0.9)$ m (the root zone plus a margin), so that the allowed head is $h=z_{drain}-z_c$.

\begin{table}[!htbp]\centering\footnotesize\setlength{\tabcolsep}{3pt}\caption{Hooghoudt spacing for the loam S06 ($K_1=K_2=0.29$ m/d, $q=6$ mm/d, $z_c=0.5$ m) at several drain depths and barrier depths.}
\begin{tabular}{rrrrrr}\toprule
\hd{Drain depth (m)} & \hd{Barrier $D$ (m)} & \hd{$h$ (m)} & \hd{$d_e$ (m)} & \hd{$L$ (m)} & \hd{pipe (m/ha)}\\\midrule
1.0 & 2 & 0.50 & 1.11 & 16.2 & 625\\
1.0 & 4 & 0.50 & 1.36 & 17.9 & 588\\
1.0 & 8 & 0.50 & 1.49 & 18.5 & 556\\
1.2 & 2 & 0.70 & 1.21 & 20.7 & 500\\
1.2 & 4 & 0.70 & 1.64 & 23.3 & 435\\
1.2 & 8 & 0.70 & 1.85 & 24.6 & 417\\
1.5 & 2 & 1.00 & 1.33 & 26.7 & 385\\
1.5 & 4 & 1.00 & 1.91 & 30.6 & 333\\
1.5 & 8 & 1.00 & 2.34 & 33.1 & 303\\
1.8 & 2 & 1.30 & 1.42 & 32.3 & 312\\
1.8 & 4 & 1.30 & 2.12 & 37.3 & 270\\
1.8 & 8 & 1.30 & 2.71 & 41.1 & 244\\
2.0 & 2 & 1.50 & 1.46 & 35.9 & 286\\
2.0 & 4 & 1.50 & 2.22 & 41.6 & 244\\
2.0 & 8 & 1.50 & 2.92 & 46.1 & 217\\\bottomrule\end{tabular}\end{table}

\exbox{Drain spacing by Hooghoudt}{Loam S06 ($K_1=K_2=0.290$ m\,d$^{-1}$, $\mu=0.164$). Drains of 100 mm lie at 1.2 m with the barrier $D=3$ m below them; the design recharge is 6 mm\,d$^{-1}$ and the water table must not rise within 0.5 m of the surface. Find the allowed head, the spacing and the water table at that spacing.}\sol $h=1.2-0.5=0.70$ m; $u=\pi(0.05)=0.157$ m \ans First estimate: $L_0=\sqrt{(4K_1h^2+8K_2Dh)/q}=\sqrt{(4\times0.290\times0.49+8\times0.290\times3\times0.7)/0.006}=30.1$ m \ans Equivalent depth at $L_0$ (Eq.~\ref{eq:dr-de}): $x=2\pi(3)/30.1$, $d_e=1.72$ m; new $L=\sqrt{(8K_2d_eh+4K_1h^2)/q}$ \ans After iteration $L=22.3$ m, $d_e=1.48$ m; the spacing is taken as $\lfloor L\rfloor=22$ m \ans Check: $h_{mid}=0.687$ m at that spacing from the quadratic $4K_1h^2+8K_2d_eh-qL^2=0$ \ans Pipe per hectare $10^4/22=455$ m \ans

\section{Transient design: Glover--Dumm}
For a rain that raises the water table, the farmer wants to know how fast it falls. The \emph{falling-water-table} basis of the application gives the spacing that lowers the water table from $h_0=z_{drain}-0.1$ to the allowed height $h_t$ in $t$ days:

\begin{equation}\boxed{L=\pi\sqrt{\frac{K\,d_e\,t}{\mu\,\ln\bigl(1.16\,h_0/h_t\bigr)}},\qquad K=\frac{K_1\bar h+K_2d_e}{\bar h+d_e},\quad\bar h=\tfrac12(h_0+h_t),}\label{eq:dr-gd}\end{equation}
\par\noindent where $L$ is the spacing (m), $K$ the weighted conductivity (m\,d$^{-1}$), $d_e$ the equivalent depth (m), $t$ the time allowed to draw the water table down (d; default 5), $\mu$ the drainable porosity ($\theta_s-\theta_{33}$, between 0.03 and 0.25), $h_0$ and $h_t$ the initial and final heights of the water table above the drain level (m) and 1.16 is the factor of the Glover--Dumm solution for a water table that starts as a flat horizontal line.



The design basis may be the steady spacing, the falling spacing, or the smaller of the two (the default for a drain that must satisfy both). The water table midway for any chosen $L$ follows from the quadratic of Hooghoudt, $h_{mid}=\bigl[-8K_2d_e+\sqrt{(8K_2d_e)^2+16K_1qL^2}\bigr]/(8K_1)$, and the discharge for any height $h$ is $q=(8K_2d_eh+4K_1h^2)/L^2$.

\begin{table}[!htbp]\centering\footnotesize\setlength{\tabcolsep}{3pt}\caption{Steady and falling-water-table spacings for the loam S06 ($D=3$ m, drain at 1.2 m, $q=6$ mm/d, $z_c=0.5$ m).}
\begin{tabular}{rrrr}\toprule
\hd{draw-down time $t$ (d)} & \hd{steady $L$ (m)} & \hd{falling $L$ (m)} & \hd{design (smaller) (m)}\\\midrule
2 & 22.3 & 6.2 & 6.2\\
3 & 22.3 & 8.5 & 8.5\\
5 & 22.3 & 12.4 & 12.4\\
7 & 22.3 & 15.8 & 15.8\\
10 & 22.3 & 20.3 & 20.3\\\bottomrule\end{tabular}\end{table}

\section{The ten soils}
\begin{table}[!htbp]\centering\footnotesize\setlength{\tabcolsep}{3pt}\caption{Drain spacing for each of the ten soils (drain 1.2 m, barrier 3 m, $q=6$ mm/d, $z_c=0.5$ m, $t=5$ d); $L$ is limited to 5--200 m.}
\begin{tabular}{llrrrrrr}\toprule
\hd{ID} & \hd{Soil} & \hd{$K_1$ (m/d)} & \hd{$\mu$} & \hd{steady $L$ (m)} & \hd{falling $L$ (m)} & \hd{$L$ used (m)} & \hd{pipe (m/ha)}\\\midrule
S01 & Coarse loam & 0.316 & 0.171 & 23.5 & 12.9 & 23 & 435\\
S02 & Coarse sand & 2.754 & 0.250 & 83.5 & 42.1 & 83 & 120\\
S03 & Sand & 3.233 & 0.250 & 91.2 & 46.3 & 91 & 110\\
S04 & Loamy sand & 2.073 & 0.250 & 71.3 & 35.3 & 71 & 141\\
S05 & Sandy loam & 1.073 & 0.250 & 49.1 & 23.2 & 49 & 204\\
S06 & Loam & 0.290 & 0.164 & 22.3 & 12.4 & 22 & 455\\
S07 & Silt loam & 0.262 & 0.150 & 21.0 & 12.4 & 20 & 500\\
S08 & Sandy clay loam & 0.127 & 0.129 & 13.1 & 8.0 & 13 & 769\\
S09 & Clay loam & 0.077 & 0.108 & 9.5 & 6.3 & 9 & 1111\\
S10 & Silty clay loam & 0.090 & 0.112 & 10.5 & 6.9 & 10 & 1000\\\bottomrule\end{tabular}\end{table}

The conductivity of the soil governs everything: on the sands a single drain serves 200 m of field or more (the limit of the application) and drainage is rarely necessary; on the silty and clay soils the spacing falls to a few metres, the pipe per hectare multiplies, and drainage becomes expensive. This is one reason why the drainage layer is read together with the irrigation design: a design with less deep percolation lowers the recharge $q$ and widens the spacing.

\section{Pipes and structures}
The lateral pipe must carry the discharge collected from a strip of width $L$ along its length $L_f$, $Q_l=q\,L\,L_f/86400$ (m$^3$\,s$^{-1}$, $q$ in m\,d$^{-1}$), and the collector the discharge of the whole field, $Q_c=qL_fW_f/86400$. The smallest standard diameter from 80, 100, 125, 160, 200, 250, 315, 400, 500 mm that carries the flow at 80 per cent of the full-pipe capacity is chosen. For a full circular pipe, Manning's equation reduces to

\begin{equation}\boxed{Q_{cap}=\frac1n\,\frac{\pi D^2}{4}\left(\frac D4\right)^{2/3}s^{1/2},\qquad L_{\max}=\min\left\{1500,\ \frac{0.8\,Q_{cap}}{q\,L/86400}\right\},}\label{eq:dr-pipe}\end{equation}
\par\noindent where $Q_{cap}$ is the full-pipe capacity (m$^3$\,s$^{-1}$), $n$ the Manning coefficient (0.016 for the pipe), $D$ the inside diameter (m), $s$ the slope of the pipe (m\,m$^{-1}$), $L_{\max}$ the longest lateral that the pipe can serve (m, limited to 1500 m), $q$ the design discharge (m\,d$^{-1}$) and $L$ the drain spacing (m).



\begin{table}[!htbp]\centering\footnotesize\setlength{\tabcolsep}{3pt}\caption{Full-pipe capacity of corrugated pipe ($n=0.016$) at three slopes.}
\begin{tabular}{rrrr}\toprule
\hd{$D$ (mm)} & \hd{$Q_{cap}$ at 0.15\% (L/s)} & \hd{at 0.30\% (L/s)} & \hd{at 0.50\% (L/s)}\\\midrule
80 & 0.9 & 1.3 & 1.6\\
100 & 1.6 & 2.3 & 3.0\\
125 & 2.9 & 4.2 & 5.4\\
160 & 5.7 & 8.1 & 10.4\\
200 & 10.3 & 14.6 & 18.8\\
250 & 18.7 & 26.5 & 34.2\\
315 & 34.7 & 49.0 & 63.3\\
400 & 65.5 & 92.7 & 119.6\\
500 & 118.8 & 168.0 & 216.9\\\bottomrule\end{tabular}\end{table}

\exbox{Pipes of a drained field}{A field of 300 m by 200 m on the loam is drained with laterals along the 300-m side, at the spacing of the previous example, slope 0.15 per cent, $q=6$ mm\,d$^{-1}$. Find the number of laterals, the flow in a lateral, the pipe size, the longest lateral and the collector.}\sol Spacing $L=22$ m; number of laterals $\lfloor200/22\rfloor+1=10$ \ans $Q_l=qLL_f/86400=0.006\times22\times300/86400=4.58e-04$ m$^3$\,s$^{-1}$ ($0.458$ L\,s$^{-1}$) \ans Smallest pipe with $0.8Q_{cap}\ge Q_l$ at $s=0.0015$: $D=80$ mm ($Q_{cap}=0.90$ L\,s$^{-1}$) \ans $L_{\max}=\min\{1500,\ 0.8Q_{cap}/(qL/86400)\}=469$ m \ans Collector: $Q_c=0.006\times300\times200/86400=4.17e-03$ m$^3$\,s$^{-1}$ ($4.17$ L\,s$^{-1}$), diameter $160$ mm \ans Pipe in the field: $10\times300+200=3200$ m, i.e.\ $533$ m\,ha$^{-1}$ over $6.0$ ha \ans

\section{Salinity and the leaching requirement}
Drainage and salinity are one problem. Each irrigation brings salt; the crop takes the water and leaves the salt; unless a fraction of the water passes below the roots the salt accumulates. Rhoades' leaching requirement gives the fraction of the water that must leave the root zone for the crop to see no more than its tolerable salinity,

\begin{equation}\boxed{LR=\mathrm{clip}\left\{\frac{EC_w}{5EC_e-EC_w},\ 0,\ 0.9\right\},\qquad I_{leach}=\frac{LR}{1-LR}\,I_n,}\label{eq:dr-lr}\end{equation}
\par\noindent where $LR$ is the leaching requirement (fraction of the applied water), $EC_w$ the electrical conductivity of the irrigation water (dS\,m$^{-1}$), $EC_e$ the salinity of the saturation extract that the crop tolerates without loss (dS\,m$^{-1}$; the user enters the threshold of the crop), $I_{leach}$ the extra water to be applied for leaching (mm) and $I_n$ the net irrigation of the season (mm).



The deep percolation caused by the irrigation (the recharge) already provides some leaching; the design is complete when the leaching fraction of the field water balance is at least $LR$ and the drains carry that water away.

\begin{table}[!htbp]\centering\footnotesize\setlength{\tabcolsep}{3pt}\caption{Leaching requirement for crops of different tolerance and waters of different salinity.}
\begin{tabular}{rrrrrr}\toprule
\hd{$EC_e$ (dS/m)} & \hd{$EC_w=0.5$} & \hd{1.0} & \hd{1.5} & \hd{2.5} & \hd{4.0}\\\midrule
1.2 & 0.091 & 0.200 & 0.333 & 0.714 & 0.900\\
1.7 & 0.062 & 0.133 & 0.214 & 0.417 & 0.889\\
3.0 & 0.034 & 0.071 & 0.111 & 0.200 & 0.364\\
4.8 & 0.021 & 0.043 & 0.067 & 0.116 & 0.200\\
7.0 & 0.014 & 0.029 & 0.045 & 0.077 & 0.129\\
7.7 & 0.013 & 0.027 & 0.041 & 0.069 & 0.116\\\bottomrule\end{tabular}\end{table}

\exbox{Leaching for a sensitive and a tolerant crop}{Maize ($EC_e$ threshold 1.7 dS\,m$^{-1}$) is irrigated with water of 1.0 dS\,m$^{-1}$; the net irrigation of the season is $I_n=480$ mm. Cotton ($EC_e=7.7$) is irrigated with the same water. Find the leaching requirement and the extra water in each case.}\sol Maize: $LR=1.0/(5\times1.7-1.0)=0.133$ \ans $I_{leach}=\frac{0.133}{1-0.133}\times480=74$ mm \ans Cotton: $LR=1.0/(5\times7.7-1.0)=0.027$ \ans $I_{leach}=\frac{0.027}{1-0.027}\times480=13$ mm \ans The tolerant crop needs one-fifth of the leaching of the sensitive one; with water of 1 dS\,m$^{-1}$ the maize needs about 15 per cent more water, which the drains must carry away \ans

\section{Controlled drainage}
With a gate on the outlet the drain does not discharge until the water table rises above the gate level $g$. The water table then follows, day by day, the balance

\begin{equation}\boxed{\mu\frac{dh}{dt}=R(t)-q_{out}(h-g),\qquad q_{out}(h')=\frac{8K_2d_eh'+4K_1h'^2}{L^2}\ (h'>0),}\label{eq:dr-ctl}\end{equation}
\par\noindent where $\mu$ is the drainable porosity, $h$ the height of the water table midway above the drain level (m), $g$ the height of the gate (m), $R(t)$ the recharge of the day (m\,d$^{-1}$), $q_{out}$ the discharge of the drain (m\,d$^{-1}$) and $h'=h-g$ the head above the gate. The application integrates it in six sub-steps per day over the whole record, with the water table capped at the soil surface (the excess is runoff) and reports the days on which the water table is shallower than the critical depth $z_c$, the sum of the exceedance of the depth of 0.3 m (the SEW$_{30}$ index, in cm\,d), and the total drainage.



Raising the gate from 0 to 0.3 m keeps more of the winter rain in the profile and cuts the discharge (and the loss of nitrate with it), at the price of more days with a shallow water table. The balance between these two is what the controlled-drainage design is for.

}{}
\banner{Drainage design}\label{ch:drain}
\section{The depth to be kept free}
The crop sets how deep the water table must stay: $z_c=0.3+0.4R_{\max}$ (m), bounded to 0.4--0.9 m, with $R_{\max}$ the rooting depth. The allowed head of the water table above the drain level, midway between the drains, is then $h=z_d-z_c$ (drain depth less $z_c$).

\exbox{The water table limit}{Maize roots to 1.0 m; drains lie 1.4 m deep. Find $z_c$ and $h$. Then repeat for a shallow-rooted crop ($R_{\max}=0.3$ m) and a deep one ($R_{\max}=1.8$ m).}\sol \givenline{$R_{\max}=0.3$ m ; $R_{\max}=1.8$ m}\ans Maize: $z_c=0.3+0.4\times1.0=0.70$ m, within 0.4--0.9, so $h=1.4-0.70=0.70$ m. Shallow crop: $0.3+0.12=0.42$ m, so $h=0.98$ m: the drains may let the water table rise higher between them. Deep crop: $0.3+0.72=1.02$ m, clipped to 0.9 m, so $h=0.50$ m. The deeper the roots, the closer the drains must be. \finans


\section{Recharge and the design discharge}
The recharge is the deep percolation of the daily balance plus the part of the applied irrigation lost to deep percolation by the irrigation design in use (DP times gross water, gross $=$ net$/E_a$). The design discharge $q$ is the 95th percentile of the three-day mean recharge of the season, between 2 and 25 mm d$^{-1}$.

\exbox{Design discharge from the daily balance}{The daily bucket model (Chapter~\ref{ch:e2}) of the loam is run for a wet year ($P=900$ mm, seed 3). It gives rainfall 808 mm, irrigation 105 mm, deep percolation 386 mm. The sprinkler layout of Chapter~\ref{ch:sprinkler} loses $\mathrm{DP}=6.0\%$ of the gross, $E_a=82\%$, so each irrigation adds $0.060/0.820=0.0732$ times its net depth to the recharge: $7.7$ mm over the season. Total recharge $=386+7.7=394$ mm. The three-day means have 95th percentile $6.43$ mm d$^{-1}$, inside 2--25 mm d$^{-1}$, so $q=6.43$ mm d$^{-1}=0.00643$ m d$^{-1}$.}
\textbf{Reading.} By construction 5\% of the three-day windows carry more recharge than the design value; the drains are not sized for those.


\section{Spacing: Hooghoudt with the exact equivalent depth}
For steady flow, with recharge $q$ (m d$^{-1}$), water table height $h$ midway, conductivities $K_1$ above and $K_2$ below the drain, the spacing $L$ is
\begin{equation}\boxed{qL^2=8K_2dh+4K_1h^2.}\end{equation}
\par\noindent where $q$ is the drain discharge per unit area (m\,d$^{-1}$), $L$ the drain spacing (m), $K_1$ and $K_2$ the conductivities above and below the drains (m\,d$^{-1}$), $d$ the equivalent depth of the barrier below the drains (m) and $h$ the height of the water table midway between drains above the drain level (m).

The equivalent depth $d$ replaces the real depth $D$ of the barrier below the drain because the flow converges radially on the drain; the layer solves the series of Moody and van Beers rather than reading a table:
\begin{equation}\boxed{d=\frac{\pi L/8}{\ln(L/u)+F(x)},\quad F(x)=\sum_{n=1,3,5,\dots}\frac{4e^{-2nx}}{n(1-e^{-2nx})},\quad x=\frac{2\pi D}L,\quad d\le D,}\end{equation}
with $u=\pi r_0$ the wetted perimeter of a pipe running half full. Because $d$ depends on $L$, the pair is iterated.

\begin{table}[!htbp]\centering\small\caption{Iteration of the spacing: $q=6.43$ mm d$^{-1}$, $K_1=K_2=0.290$ m d$^{-1}$, $D=5$ m, $h=0.70$ m, $r_0=0.05$ m.}\label{tab:hoogit}
\begin{tabular}{lrrr}\toprule
\hd{Iteration} & \hd{$L$ used (m)} & \hd{$d$ (m)} & \hd{$L$ new (m)}\\\midrule
1 & 30.000 & 2.025 & 24.475\\
2 & 24.475 & 1.786 & 23.208\\
3 & 23.208 & 1.726 & 22.878\\
4 & 22.878 & 1.709 & 22.789\\
5 & 22.789 & 1.705 & 22.765\\
6 & 22.765 & 1.704 & 22.758\\
7 & 22.758 & 1.704 & 22.756\\
8 & 22.756 & 1.703 & 22.756\\\bottomrule\end{tabular}\end{table}

\exbox{One converged step}{Verify the converged spacing of Table~\ref{tab:hoogit}.}\sol  At $L=22.76$ m: $x=2\pi\times5/22.76=1.381$. $F(x)=4e^{-2x}/(1-e^{-2x})+\dots=0.2703$ (the $n=1$ term dominates). $u=\pi\times0.05=0.1571$ m, $\ln(L/u)=\ln(22.76/0.1571)=4.976$. Then $d=(\pi\times22.76/8)/(4.976+0.2703)=8.936/5.246=1.703$ m ($<D$). Check the spacing: $8K_2dh=8\times0.290\times1.703\times0.70=2.7631$, $4K_1h^2=4\times0.290\times0.49=0.5677$; sum $=3.3309$; $L^2=3.3309/0.00643=517.8$, $L=22.76$ m. ({\it Without} the correction, using $d=D=5$ m, the spacing would be $\sqrt{(8\times0.290\times5\times0.7+0.5677)/0.00643}=36.7$ m: the radial-flow correction matters.) \finans


\begin{table}[!htbp]\centering\small\caption{The same drains in the ten soils: $q=6.43$ mm d$^{-1}$, $z_d=1.4$ m, $h=0.70$ m, $D=5$ m, $K_2=K_1$ (Saxton--Rawls $K_s$).}\label{tab:drsoils}
\begin{tabular}{lrrrr}\toprule
\hd{Soil} & \hd{$K$ (m/d)} & \hd{$L$ (m)} & \hd{$d$ (m)} & \hd{pipe (m/ha)}\\\midrule
Coarse loam & 0.316 & 24.1 & 1.77 & 414\\
Coarse sand & 2.754 & 95.0 & 3.42 & 105\\
Sand & 3.233 & 104.3 & 3.52 & 96\\
Loamy sand & 2.073 & 80.4 & 3.23 & 124\\
Sandy loam & 1.073 & 53.8 & 2.75 & 186\\
Loam & 0.290 & 22.8 & 1.70 & 439\\
Silt loam & 0.262 & 21.2 & 1.63 & 471\\
Sandy clay loam & 0.127 & 12.8 & 1.14 & 781\\
Clay loam & 0.077 & 9.1 & 0.88 & 1{,}100\\
Silty clay loam & 0.090 & 10.1 & 0.95 & 990\\\bottomrule\end{tabular}\end{table}

\begin{keybox}Table~\ref{tab:drsoils} runs from 9.1 m (Clay loam) to 104.3 m (Sand). The spacing scales roughly with $\sqrt K$ (a tenfold increase in $K$ gives a spacing about three times larger), so the pipe per hectare, $10^4/L$ metres, is the cost driver of a drainage scheme in the finer soils.\end{keybox}
\exbox{Falling water table}{After a wet spell the water table must fall from $h_0=0.70$ m to $h_t=0.35$ m above the drain within $t=3$ d. The drainable porosity of the loam is $\mu=0.95\theta_s-\theta_{33}=0.143$ (between 0.03 and 0.25). Find the spacing from $L^2=\pi^2Kdt/[\mu\ln(1.16h_0/h_t)]$ and compare with the steady value.}\sol \givenline{$h_0=0.70$ m ; $h_t=0.35$ m above ; $t=3$ d ; $\mu=0.95\theta_s-\theta_{33}=0.143$ ; $L^2=\pi^2Kdt/[\mu\ln(1.16h_0/h_t)]$}\ans $\ln(1.16\times0.70/0.35)=\ln2.32=0.842$. $L^2=\pi^2\times0.290\times1.703\times3/(0.143\times0.842)=14.609/0.1202=121.5$, $L=11.0$ m. The steady (Hooghoudt) spacing is 22.8 m, so the design takes the closer, $L=\min=11.0$ m. (Using $d=1.703$ m from the steady solution; the layer iterates $d$ with $L$ in the same way.) \finans


\subhd{How to read the spacing equation.} The two terms of Hooghoudt's equation have a plain meaning. Water that is added to the soil by rain or irrigation reaches the water table, and the water table rises until the water can leave towards the drain at the same rate as it arrives. The flow towards the drain has two parts. One part moves through the soil \emph{below} the level of the drain, in a thick layer of depth $d$, and gives the term $8K_2dh$. The other part moves \emph{above} the level of the drain, in the layer between the drain level and the water table, and gives the term $4K_1h^2$. A larger spacing $L$ means that the water must travel further, so the water table between the drains rises higher. Closer drains keep the water table lower. The designer's task is to choose the largest $L$ for which the water table midway between two drains does not rise above the allowed height $h$.

\exbox{One pass of the spacing calculation by hand}{Use the data of Table~\ref{tab:hoogit}: $q=6.43$ mm d$^{-1}$, $K_1=K_2=0.290$ m d$^{-1}$, allowed height $h=0.70$ m. (a) With the first value of the equivalent depth, $d=2.025$ m (found for a trial spacing of 30 m), compute the spacing. (b) Repeat with the improved depth $d=1.786$ m. (c) Explain why the calculation must be repeated.}\sol \givenline{$q=6.43\times10^{-3}$ m d$^{-1}$ ; $K_1=K_2=0.290$ m d$^{-1}$ ; $h=0.70$ m}
\par\noindent We know that $qL^2=8K_2dh+4K_1h^2$, so $L=\sqrt{(8K_2dh+4K_1h^2)/q}$.
\par\noindent The second term does not depend on the spacing: $4K_1h^2=4\times0.290\times0.70^2=0.568$ m$^2$ d$^{-1}$.
\par\noindent\textbf{(a) First pass.} With $d=2.025$ m,
\[8K_2dh=8\times0.290\times2.025\times0.70=3.289\ \text{m}^2\,\text{d}^{-1},\]
\[L=\sqrt{\frac{3.289+0.568}{0.00643}}=\sqrt{599.8}=24.5\ \text{m}\quad\text{\finans}\]
\par\noindent\textbf{(b) Second pass.} With $d=1.786$ m,
\[8K_2dh=8\times0.290\times1.786\times0.70=2.900,\qquad L=\sqrt{\frac{2.900+0.568}{0.00643}}=\sqrt{539.4}=23.2\ \text{m}\quad\text{\finans}\]
\par\noindent\textbf{(c) Why repeat.} The equivalent depth $d$ is not fixed: it depends on the spacing itself, because the closer the drains, the more sharply the flow bends towards each of them. A trial spacing gives a depth $d$, which gives a new spacing, which gives a new depth. Table~\ref{tab:hoogit} shows that after eight such passes the two agree at $L=22.76$ m. The first pass is already within 8 per cent of the final answer, and the third within 1 per cent. In practice the drains are laid at the nearest convenient spacing below this value, here $22$ m.

\section{The water table through the season}
The daily balance is $\mu\,dh/dt=R-q_{\rm out}(h-g)$ with the outflow $q_{\rm out}(h)=(8K_2dh+4K_1h^2)/L^2$ (six sub-steps a day), $R$ the recharge of the day and $g$ the height of the control gate (zero for plain drains). The results are the days with the water table shallower than the crop needs, the $\mathrm{SEW}_{30}$ of Sieben (the sum over the days of $30-z$, in cm, when the water table is shallower than 30 cm), the drained water, the maximum outflow, and for the gate the water held, $\mu g$.

\begin{table}[!htbp]\centering\small\caption{Water-table results for the wet-year recharge at four spacings ($d$ kept at the design value).}\label{tab:wtsim}
\begin{tabular}{lrrrr}\toprule
\hd{$L$ (m)} & \hd{days $z<z_c$} & \hd{SEW$_{30}$} & \hd{drained (mm)} & \hd{max outflow (mm/d)}\\\midrule
11.0 & 0 & 0 & 371 & 8.02\\
22.8 & 0 & 0 & 343 & 4.25\\
35.0 & 8 & 0 & 296 & 3.16\\
50.0 & 59 & 3 & 242 & 2.32\\
70.0 & 68 & 710 & 181 & 1.59\\\bottomrule\end{tabular}\end{table}

\begin{table}[!htbp]\centering\small\caption{Controlled drainage: gate height at the design spacing.}\label{tab:gate}
\begin{tabular}{lrrrr}\toprule
\hd{gate $g$ (m)} & \hd{days $z<z_c$} & \hd{SEW$_{30}$} & \hd{drained (mm)} & \hd{water held $\mu g$ (mm)}\\\midrule
0.00 & 0 & 0 & 371 & 0.0\\
0.10 & 0 & 0 & 356 & 14.3\\
0.20 & 0 & 0 & 342 & 28.6\\
0.30 & 0 & 0 & 328 & 42.8\\\bottomrule\end{tabular}\end{table}

\exbox{Spacing and gate in the season}{Table~\ref{tab:wtsim} and Table~\ref{tab:gate} come from the same recharge series. Read what they say about spacing and the gate.}\sol  At the design spacing $L=11.0$ m the water table is shallower than $z_c=0.70$ m on 0 day(s). Widening the spacing to 50 m raises that to 59 days (SEW$_{30}$ = 3) and to 70 m, 68 days. A gate raised by 0.3 m holds $\mu g=43$ mm in the soil (the loam's $\mu=0.143$) and changes the number of days shallower than $z_c$ to 0. The gate saves water for the crop but costs aeration if it is raised beyond the largest height that keeps the water table below the crop; the layer displays that limit. \finans


\section{Pipes}
Drain pipes are sized by the Manning equation for a full pipe,
\begin{equation}\boxed{Q=\frac1n\,\frac{\pi D^2}4\left(\frac D4\right)^{2/3}s^{1/2},\quad n=0.016\ \text{(corrugated polyethylene)},}\end{equation}
\par\noindent where $Q$ is the full-pipe capacity (m$^3$\,s$^{-1}$), $n$ Manning's roughness coefficient, $D$ the pipe diameter (m) and $s$ the slope (m\,m$^{-1}$); $\pi D^2/4$ is the area and $D/4$ the hydraulic radius of a full circular pipe.

choosing the smallest standard size (80--500 mm) that carries the design flow at no more than 80\% of its capacity. The lateral flow is $qLL_{\rm lat}$ and the collector flow is $qA$. Silt settles when the full-pipe velocity is below 0.3 m s$^{-1}$.

\begin{table}[!htbp]\centering\small\caption{Manning capacity of corrugated pipe at slope 0.2\%.}\label{tab:pipe}
\begin{tabular}{lrrr}\toprule
\hd{$D$ (mm)} & \hd{$Q_{\rm full}$ (L/s)} & \hd{$0.8Q$ (L/s)} & \hd{$v$ (m/s)}\\\midrule
80 & 1.04 & 0.83 & 0.21\\
100 & 1.88 & 1.50 & 0.24\\
125 & 3.40 & 2.72 & 0.28\\
160 & 6.57 & 5.26 & 0.33\\
200 & 11.92 & 9.53 & 0.38\\
250 & 21.61 & 17.29 & 0.44\\
315 & 40.02 & 32.02 & 0.51\\
400 & 75.67 & 60.54 & 0.60\\
500 & 137.20 & 109.76 & 0.70\\\bottomrule\end{tabular}\end{table}

\exbox{Pipes of the scheme}{The field is 400 m $\times$ 300 m; laterals run 200 m at slope 0.2\%. Size the lateral and the collector.}\sol  The lateral drains a strip $L=11.0$ m wide and 200 m long: $Q=qLL_{\rm lat}=0.00643\times11.0\times200=14.18$ m$^3$ d$^{-1}=0.164$ L s$^{-1}$. With 80\% of capacity allowed, Table~\ref{tab:pipe} gives the smallest pipe whose $0.8Q$ exceeds 0.164 L s$^{-1}$: $D=80$ mm. The collector carries the whole 12 ha: $Q=qA=0.00643\times120{,}000=771.9$ m$^3$ d$^{-1}=8.93$ L s$^{-1}$, so $D=200$ mm. The full-pipe speed of the lateral is 0.21 m s$^{-1}$, below the silting limit of 0.3 m s$^{-1}$. \finans


\section{Salt and the leaching requirement}
The leaching requirement of Rhoades is $LR=EC_w/(5EC_e-EC_w)$, with $EC_w$ the salinity of the irrigation water and $EC_e$ the salinity of the saturation extract the crop tolerates (dS m$^{-1}$); the gross irrigation is $\mathrm{net}/(1-LR)$ when $LR$ is added to the need.

\begin{table}[!htbp]\centering\small\caption{Leaching requirement for two water qualities.}\label{tab:lr}
\begin{tabular}{lrrrr}\toprule
\hd{$EC_e$} & \hd{$EC_w$} & \hd{$LR$} & \hd{$EC_w$} & \hd{$LR$}\\\midrule
1.5 & 0.5 & 0.0714 & 1.5 & 0.2500\\
3 & 0.5 & 0.0345 & 1.5 & 0.1111\\
6 & 0.5 & 0.0169 & 1.5 & 0.0526\\
10 & 0.5 & 0.0101 & 1.5 & 0.0309\\\bottomrule\end{tabular}\end{table}

\exbox{Leaching requirement}{Irrigation water of 1.5 dS m$^{-1}$ is applied to a crop that tolerates $EC_e=3$ dS m$^{-1}$ (a moderately sensitive crop), net need 211.68 mm. Find $LR$ and the gross.}\sol \givenline{$EC_e=3$ dS m$^{-1}$}\ans $LR=1.5/(5\times3-1.5)=1.5/13.5=0.1111$. Gross $=211.68/(1-0.1111)=238.1$ mm: an extra 26.5 mm, which must leave through the drains. For $EC_e=6$: $LR=1.5/28.5=0.0526$ and the extra is only 11.8 mm. The more tolerant the crop, the less water the drains have to carry away for salt. \finans


\section{Summary of the drainage layer}
\begin{itemize}
\item Depth limit $z_c=0.3+0.4R_{\max}$; allowed head $h=z_d-z_c$.
\item Design discharge $=$ 95th percentile of the three-day recharge (rain and irrigation losses).
\item Spacing: Hooghoudt with Moody--van Beers equivalent depth; Glover--Dumm for the falling table; the closer governs.
\item Season: $\mu\,dh/dt=R-q_{\rm out}(h-g)$; $\mathrm{SEW}_{30}$; gate holds $\mu g$.
\item Pipes by Manning at 80\% capacity; leaching $LR=EC_w/(5EC_e-EC_w)$.
\end{itemize}
\IfFileExists{extra/app14.tex}{\subsection{The output of the application, panel by panel}
The drainage layer is a numbered sequence of six panels that answers three questions: how much water must the drains remove, how far apart must they be, and how large must the pipes be. The case shown is a buried pipe drain at 1.20 m in the coarse loam S01, with a barrier 3.0 m below the drains, a field of 300 by 200 m and a pipe slope of 0.15 per cent, designed by the steady-flow (Hooghoudt) method with the discharge taken from the DSS recharge.

The buttons choose the type of drain and the design basis of this chapter; spacing and design flow can be automatic (from the equations and the DSS recharge) or typed by the user. The boxes hold the drain depth (1.20 m), the barrier $D$ (3.0 m), the subsoil conductivity as a multiple of that of the topsoil, the time to lower the water table (5.0 d), the field and lateral length, the pipe slope, the critical depth of the water table (0.68 m, from the roots of the crop), the gate height, the ditch width, the salt tolerance of the crop and the lift to the outlet.

\appfig{drn_inputs}{0.96}{Inputs of the drainage layer: crop, mulch, drain type (pipe, pipe with control gate, open ditch), design basis (steady flow, falling water table, both), spacing and flow modes, and the numerical variables.}

The green box states the design as a recommendation a farmer can act on, with the check that the water table stays 0.69 m below the surface, deeper than the 0.68 m the roots need, the number of days too wet in the season (0 of 121), the water carried away by the drains (48 mm), and the reasoning in a sentence.

\appfig{drn_answer}{0.96}{The answer in plain words: one pipe drain at 1.20 m every 37 m, 6 pipes of 300 m with 80 mm pipe and a 125 mm collector.}

The bars are the daily recharge $q$ in mm\,d$^{-1}$ and the dashed line is the design discharge, the value exceeded on only 5 per cent of days (2.0 mm\,d$^{-1}$, the $P_{95}$ of this chapter). The table gives the rain of the season (176 mm), the net irrigation, the efficiency of the irrigation design (43 per cent) and the share of the applied water that percolates (7 per cent).

\appfig{drn_recharge}{0.96}{1. How much water must the drains remove? Day-by-day deep percolation of the DSS season, from rain and from the part of the irrigation that is lost.}

The curves are the Hooghoudt equation solved for the spacing: the lower the discharge the wider the spacing, and a deeper drain (green) allows a wider spacing than a shallow one (brown). The red dot is the design ($q=2.0$ mm\,d$^{-1}$, $L=37$ m). The table lists the conductivity above and below the drain (0.32 m\,d$^{-1}$), the drainable pore space (0.193), the effective depth $d=1.86$ m, the rise allowed $h=0.52$ m and the spacings by the two methods (37 m steady, 9 m falling-table).

\appfig{drn_spacing}{0.96}{2. How far apart must the drains be? The spacing $L$ against the design discharge $q$ for three drain depths.}

The picture shows the water-filled soil, the critical depth of 0.68 m (red dashed line) and the water table midway between the drains at 1.15 m. Both pictures are the same until a control gate is chosen, when the water table on the right is held higher.

\appfig{drn_side}{0.96}{3. Side view between two drains, drawn to scale: free drainage (left) and the pipe with a gate (right).}

The upper chart follows the depth of the water table through the 121 days against the drain level and the critical depth; the lower chart is the outflow of the drains in mm\,d$^{-1}$ against the design value. The table lists the days too wet, the SEW$_{30}$ index, the water drained, the largest outflow, the runoff, the leaching requirement and the water needed for salt.

\appfig{drn_wt}{0.96}{4. Where the water table sits through the season, and the outflow of the drains.}

The plan shows the six laterals of 300 m with the collector and the ditch; the table lists the flow of a lateral (0.26 L\,s$^{-1}$), the lateral pipe (80 mm), the longest lateral that the pipe can carry (837 m), the collector flow (1.4 L\,s$^{-1}$) and its pipe (125 mm), the trench volume (1\,080 m$^3$) and the pipe laid (2\,000 m). The pipes are sized with Manning's equation ($n=0.016$, corrugated pipe) running at most 80 per cent full.

\appfig{drn_layout}{0.96}{5. Layout: pipes, collector and outlet, seen from above.}

The table repeats the calculation for the ten soils of the application at the same discharge and depth, so the effect of the soil is isolated: the sand needs a spacing of 138 m and the silty clay loam 17 m, a difference of eight times. The row of the chosen soil is highlighted.

\appfig{drn_soils}{0.96}{6. The same drains in every soil: spacing at the design discharge and depth for the ten soils, from their conductivity (Saxton and Rawls).}

The cards give the result in the units a contractor needs: the length of pipe per hectare, the drainage coefficient in mm\,d$^{-1}$ and in m$^3$ per hectare per day, the depth of the water table midway and whether it meets the critical depth, the days too wet, the volume drained and the diameters of the pipes.

\appfig{drn_result}{0.96}{Result for this design: drain spacing 37 m (333 m per hectare), drainage coefficient 2.0 mm\,d$^{-1}$, water table midway 0.69 m deep, 0 days too wet, 48 mm drained, lateral 80 mm, collector 125 mm.}

The last table puts the answer in farm terms: the field of 6.0 ha needs 6 laterals every 37 m, the drained water is 475 m$^3$ per hectare for the season, and the salt balance says the water of 0.6 dS\,m$^{-1}$ needs a 4 per cent leaching requirement, which the drained water already provides. A note warns when the flow speed in the pipe falls below 0.3 m\,s$^{-1}$, when silt would settle and a steeper slope is needed.

\appfig{drn_meaning}{0.96}{What this means for the farm: the soil, crop and irrigation in use, the water table to be kept, the water per hectare, the salt balance and the total.}

}{}
\IfFileExists{extra/long14.tex}{\subsection{Complete design: a subsurface drainage system for 20 ha}
\exbox{Design of a subsurface drainage system for 20 ha}{A field of 20 ha of loam S06 needs drainage. The pipe drains (100 mm) are laid at 1.4 m on a barrier 3 m below, the design recharge is 7 mm\,d$^{-1}$ and the water table must stay 0.6 m below the surface. Laterals are 300 m long with a grade of 0.2 per cent. Find the spacing, the number and total length of laterals, the discharge and pipe size, and the outlet flow.}\sol
\ans Soil S06 (loam): $K_s=12.1$ mm\,h$^{-1}$, so $K_1=K_2=0.290$ m\,d$^{-1}$; drainable porosity \[\mu=\theta_s-\theta_{33}=0.164\]
\ans Design recharge $q=7$ mm\,d$^{-1}$, drains at $1.4$ m on a barrier $D=3$ m below, critical water table $z_c=0.6$ m, so the allowed head is $h=1.4-0.6=0.80$ m
\ans Perimeter of the pipe $u=\pi r_0=\pi(0.05)=0.157$ m
\ans First guess for the spacing with the real depth, \[L_0=\sqrt{(4K_1h^2+8K_2Dh)/q}=\sqrt{(4\times0.290\times0.64+8\times0.290\times3\times0.80)/0.007}=30\ \text{m}\]
\ans Equivalent depth at this spacing: $x=2\pi D/L=0.628$; the series $F(x)$ and Hooghoudt's expression give $d_e=1.71$ m, much less than $D$
\ans Iterating $L=\sqrt{(8K_2d_eh+4K_1h^2)/q}$ to a change below 1 cm: $L=22.4$ m with $d_e=1.50$ m; adopt $L=22$ m
\ans Number of laterals of 300 m: $20$ ha $=200\,000$ m$^2$; $200\,000/(22\times300)=31$ laterals, total pipe $9300$ m, or $465$ m per hectare
\ans Discharge at the lower end of a lateral: \[Q=q\,L\,\ell=0.007\times22\times300=46.2\] m$^3$\,d$^{-1}$ $=0.53$ L\,s$^{-1}$
\ans Manning for a pipe running full, $n=0.016$, slope 0.2 per cent: \[Q=\dfrac{1}{n}\dfrac{\pi D_p^2}{4}\left(\dfrac{D_p}{4}\right)^{2/3}S^{1/2}\]
keeping a margin of 20 per cent the smallest standard diameter that carries $0.53$ L\,s$^{-1}$ is $80$ mm
\ans Collector: the whole area delivers \[Q_c=q\,A=0.007\times200\,000/86.4=16.2\] L\,s$^{-1}$ to the outlet
\ans Water table in the middle at this spacing, from $q=(8K_2d_eh+4K_1h^2)/L^2$: $h=0.80$ m, so the water table is 0.60 m below the surface, deeper than $z_c=0.6$ m \finans


}{}
\IfFileExists{extra/proj14.tex}{\banner{Design project 14: a drainage system for an irrigated clay field}\label{sec:proj14}
\textbf{Brief.} A 40-ha field of cotton ($EC_e$ threshold 7.7 dS\,m$^{-1}$) is irrigated with water of 1.5 dS\,m$^{-1}$ and has a barrier at 6 m. The design discharge is 7 mm\,d$^{-1}$ and the water table midway between drains is allowed to stand $h=0.5$ m above drain level; $K_1=0.8$, $K_2=0.3$ m\,d$^{-1}$. The designer chooses a drain depth, and then the spacing and the total pipe.

\begin{table}[!htbp]\centering\small\caption{Hooghoudt spacings for $q=7$ mm/d, $h=0.5$ m, $K_1=0.8$, $K_2=0.3$ m/d, barrier at 6 m.}
\begin{tabular}{rrrrr}\toprule
\hd{Drain depth (m)} & \hd{$d$ (m)} & \hd{$d_e$ (m)} & \hd{Spacing $L$ (m)} & \hd{Drain length (m/ha)}\\\midrule
1.2 & 4.8 & 1.9 & 21 & 472\\
1.5 & 4.5 & 1.9 & 21 & 475\\
1.8 & 4.2 & 1.9 & 21 & 478\\
2.1 & 3.9 & 1.8 & 21 & 482\\\bottomrule\end{tabular}\end{table}

\textbf{Reading.} Deeper drains lower the barrier-limited flow path only slightly, because the equivalent depth saturates as the drain nears the barrier, so the spacing rises little with depth; the extra excavation cost must be set against the small gain in spacing. At 1.8 m, $L=21$ m, so the field of 40 ha needs about $19$ km of pipe. The leaching requirement for cotton at an $EC_e$ of 7.7 is $0.041$, small, so the drainage is set by the water table rather than by salt.
\fig{g14_wt}{0.7}{Steady water table midway between two drains of spacing $L=21$ m (schematic elliptical shape, $h=0.5$ m).}
}{}
\IfFileExists{extra/ten14.tex}{\subsection{Worked solutions on the ten soils}
The drain spacing for a fixed discharge is computed for each of the ten soils.

\exbox{Ten-soil example 14a: Coarse loam (S01)}{Coarse loam ($K_2=0.312$ m\,d$^{-1}$ from the class value, $K_1=1.2K_2$): a drain system must remove $q=7$ mm\,d$^{-1}$ with $h=0.5$ m and an equivalent depth $d=3$ m. Find the drain spacing and the length of pipe per hectare.}\sol \givenline{$K_2=0.312$ m\,d$^{-1}$ ; $K_1=1.2K_2$ ; $q=7$ mm\,d$^{-1}$ ; $h=0.5$ m ; $d=3$ m}\ans $8K_2dh=8\times0.312\times3\times0.5=3.744$; $4K_1h^2=4\times0.374\times0.25=0.374$ \ans $L^2=(3.744+0.374)/0.007=588$, $L=24$ m \ans Pipe: $10^4/L=412$ m\,ha$^{-1}$ \ans The soil is permeable: a wide spacing is possible. \finans

\exbox{Ten-soil example 14a: Coarse sand (S02)}{Coarse sand ($K_2=2.880$ m\,d$^{-1}$ from the class value, $K_1=1.2K_2$): a drain system must remove $q=7$ mm\,d$^{-1}$ with $h=0.5$ m and an equivalent depth $d=3$ m. Find the drain spacing and the length of pipe per hectare.}\sol \givenline{$K_2=2.880$ m\,d$^{-1}$ ; $K_1=1.2K_2$ ; $q=7$ mm\,d$^{-1}$ ; $h=0.5$ m ; $d=3$ m}\ans $8K_2dh=8\times2.880\times3\times0.5=34.560$; $4K_1h^2=4\times3.456\times0.25=3.456$ \ans $L^2=(34.560+3.456)/0.007=5431$, $L=74$ m \ans Pipe: $10^4/L=136$ m\,ha$^{-1}$ \ans The soil is permeable: a wide spacing is possible. \finans

\exbox{Ten-soil example 14a: Sand (S03)}{Sand ($K_2=5.040$ m\,d$^{-1}$ from the class value, $K_1=1.2K_2$): a drain system must remove $q=7$ mm\,d$^{-1}$ with $h=0.5$ m and an equivalent depth $d=3$ m. Find the drain spacing and the length of pipe per hectare.}\sol \givenline{$K_2=5.040$ m\,d$^{-1}$ ; $K_1=1.2K_2$ ; $q=7$ mm\,d$^{-1}$ ; $h=0.5$ m ; $d=3$ m}\ans $8K_2dh=8\times5.040\times3\times0.5=60.480$; $4K_1h^2=4\times6.048\times0.25=6.048$ \ans $L^2=(60.480+6.048)/0.007=9504$, $L=97$ m \ans Pipe: $10^4/L=103$ m\,ha$^{-1}$ \ans The soil is permeable: a wide spacing is possible. \finans

\exbox{Ten-soil example 14a: Loamy sand (S04)}{Loamy sand ($K_2=1.464$ m\,d$^{-1}$ from the class value, $K_1=1.2K_2$): a drain system must remove $q=7$ mm\,d$^{-1}$ with $h=0.5$ m and an equivalent depth $d=3$ m. Find the drain spacing and the length of pipe per hectare.}\sol \givenline{$K_2=1.464$ m\,d$^{-1}$ ; $K_1=1.2K_2$ ; $q=7$ mm\,d$^{-1}$ ; $h=0.5$ m ; $d=3$ m}\ans $8K_2dh=8\times1.464\times3\times0.5=17.568$; $4K_1h^2=4\times1.757\times0.25=1.757$ \ans $L^2=(17.568+1.757)/0.007=2761$, $L=53$ m \ans Pipe: $10^4/L=190$ m\,ha$^{-1}$ \ans The soil is permeable: a wide spacing is possible. \finans

\exbox{Ten-soil example 14a: Sandy loam (S05)}{Sandy loam ($K_2=0.624$ m\,d$^{-1}$ from the class value, $K_1=1.2K_2$): a drain system must remove $q=7$ mm\,d$^{-1}$ with $h=0.5$ m and an equivalent depth $d=3$ m. Find the drain spacing and the length of pipe per hectare.}\sol \givenline{$K_2=0.624$ m\,d$^{-1}$ ; $K_1=1.2K_2$ ; $q=7$ mm\,d$^{-1}$ ; $h=0.5$ m ; $d=3$ m}\ans $8K_2dh=8\times0.624\times3\times0.5=7.488$; $4K_1h^2=4\times0.749\times0.25=0.749$ \ans $L^2=(7.488+0.749)/0.007=1177$, $L=34$ m \ans Pipe: $10^4/L=292$ m\,ha$^{-1}$ \ans The soil is permeable: a wide spacing is possible. \finans

\exbox{Ten-soil example 14a: Loam (S06)}{Loam ($K_2=0.312$ m\,d$^{-1}$ from the class value, $K_1=1.2K_2$): a drain system must remove $q=7$ mm\,d$^{-1}$ with $h=0.5$ m and an equivalent depth $d=3$ m. Find the drain spacing and the length of pipe per hectare.}\sol \givenline{$K_2=0.312$ m\,d$^{-1}$ ; $K_1=1.2K_2$ ; $q=7$ mm\,d$^{-1}$ ; $h=0.5$ m ; $d=3$ m}\ans $8K_2dh=8\times0.312\times3\times0.5=3.744$; $4K_1h^2=4\times0.374\times0.25=0.374$ \ans $L^2=(3.744+0.374)/0.007=588$, $L=24$ m \ans Pipe: $10^4/L=412$ m\,ha$^{-1}$ \ans The soil is permeable: a wide spacing is possible. \finans

\exbox{Ten-soil example 14a: Silt loam (S07)}{Silt loam ($K_2=0.163$ m\,d$^{-1}$ from the class value, $K_1=1.2K_2$): a drain system must remove $q=7$ mm\,d$^{-1}$ with $h=0.5$ m and an equivalent depth $d=3$ m. Find the drain spacing and the length of pipe per hectare.}\sol \givenline{$K_2=0.163$ m\,d$^{-1}$ ; $K_1=1.2K_2$ ; $q=7$ mm\,d$^{-1}$ ; $h=0.5$ m ; $d=3$ m}\ans $8K_2dh=8\times0.163\times3\times0.5=1.958$; $4K_1h^2=4\times0.196\times0.25=0.196$ \ans $L^2=(1.958+0.196)/0.007=308$, $L=18$ m \ans Pipe: $10^4/L=570$ m\,ha$^{-1}$ \ans The soil is slow: spacing is narrow and the cost high. \finans

\exbox{Ten-soil example 14a: Sandy clay loam (S08)}{Sandy clay loam ($K_2=0.103$ m\,d$^{-1}$ from the class value, $K_1=1.2K_2$): a drain system must remove $q=7$ mm\,d$^{-1}$ with $h=0.5$ m and an equivalent depth $d=3$ m. Find the drain spacing and the length of pipe per hectare.}\sol \givenline{$K_2=0.103$ m\,d$^{-1}$ ; $K_1=1.2K_2$ ; $q=7$ mm\,d$^{-1}$ ; $h=0.5$ m ; $d=3$ m}\ans $8K_2dh=8\times0.103\times3\times0.5=1.238$; $4K_1h^2=4\times0.124\times0.25=0.124$ \ans $L^2=(1.238+0.124)/0.007=195$, $L=14$ m \ans Pipe: $10^4/L=717$ m\,ha$^{-1}$ \ans The soil is slow: spacing is narrow and the cost high. \finans

\exbox{Ten-soil example 14a: Clay loam (S09)}{Clay loam ($K_2=0.055$ m\,d$^{-1}$ from the class value, $K_1=1.2K_2$): a drain system must remove $q=7$ mm\,d$^{-1}$ with $h=0.5$ m and an equivalent depth $d=3$ m. Find the drain spacing and the length of pipe per hectare.}\sol \givenline{$K_2=0.055$ m\,d$^{-1}$ ; $K_1=1.2K_2$ ; $q=7$ mm\,d$^{-1}$ ; $h=0.5$ m ; $d=3$ m}\ans $8K_2dh=8\times0.055\times3\times0.5=0.662$; $4K_1h^2=4\times0.066\times0.25=0.066$ \ans $L^2=(0.662+0.066)/0.007=104$, $L=10$ m \ans Pipe: $10^4/L=980$ m\,ha$^{-1}$ \ans The soil is slow: spacing is narrow and the cost high. \finans

\exbox{Ten-soil example 14a: Silty clay loam (S10)}{Silty clay loam ($K_2=0.036$ m\,d$^{-1}$ from the class value, $K_1=1.2K_2$): a drain system must remove $q=7$ mm\,d$^{-1}$ with $h=0.5$ m and an equivalent depth $d=3$ m. Find the drain spacing and the length of pipe per hectare.}\sol \givenline{$K_2=0.036$ m\,d$^{-1}$ ; $K_1=1.2K_2$ ; $q=7$ mm\,d$^{-1}$ ; $h=0.5$ m ; $d=3$ m}\ans $8K_2dh=8\times0.036\times3\times0.5=0.432$; $4K_1h^2=4\times0.043\times0.25=0.043$ \ans $L^2=(0.432+0.043)/0.007=68$, $L=8$ m \ans Pipe: $10^4/L=1214$ m\,ha$^{-1}$ \ans The soil is slow: spacing is narrow and the cost high. \finans

}{}
\IfFileExists{extra/tenb14.tex}{\subsection{Worked solutions on ten crops}
The ten crops of the library differ in their demand, their roots and their sensitivity. The same type of question is answered for each of them at $ET_0=5.5$ mm\,d$^{-1}$; the data are typical values of the crops (Doorenbos--Kassam and FAO 56).

\exbox{Crop example 14a: Maize}{Maize (threshold $EC_e=1.7$ dS\,m$^{-1}$): the irrigation water has $EC_w=1.5$ dS\,m$^{-1}$. Find the leaching requirement and the gross depth for a net requirement of 450 mm.}\sol \givenline{$EC_e=1.7$ dS\,m$^{-1}$ ; $EC_w=1.5$ dS\,m$^{-1}$}\ans \par\indent We know that
\[LR=EC_w/(5EC_e-EC_w)\ti{i}\]
Substituting the values in equation \rfi{i},
\[LR=1.5/(5\times1.7-1.5)=0.214\] \ans Gross $=450/(1-LR)=450/0.786=573$ mm \ans Extra water for leaching $123$ mm \finans

\exbox{Crop example 14a: Wheat}{Wheat (threshold $EC_e=6.0$ dS\,m$^{-1}$): the irrigation water has $EC_w=1.5$ dS\,m$^{-1}$. Find the leaching requirement and the gross depth for a net requirement of 450 mm.}\sol \givenline{$EC_e=6.0$ dS\,m$^{-1}$ ; $EC_w=1.5$ dS\,m$^{-1}$}\ans \par\indent We know that
\[LR=EC_w/(5EC_e-EC_w)\ti{i}\]
Substituting the values in equation \rfi{i},
\[LR=1.5/(5\times6.0-1.5)=0.053\] \ans Gross $=450/(1-LR)=450/0.947=475$ mm \ans Extra water for leaching $25$ mm \finans

\exbox{Crop example 14a: Cotton}{Cotton (threshold $EC_e=7.7$ dS\,m$^{-1}$): the irrigation water has $EC_w=1.5$ dS\,m$^{-1}$. Find the leaching requirement and the gross depth for a net requirement of 450 mm.}\sol \givenline{$EC_e=7.7$ dS\,m$^{-1}$ ; $EC_w=1.5$ dS\,m$^{-1}$}\ans \par\indent We know that
\[LR=EC_w/(5EC_e-EC_w)\ti{i}\]
Substituting the values in equation \rfi{i},
\[LR=1.5/(5\times7.7-1.5)=0.041\] \ans Gross $=450/(1-LR)=450/0.959=469$ mm \ans Extra water for leaching $19$ mm \finans

\exbox{Crop example 14a: Tomato}{Tomato (threshold $EC_e=2.5$ dS\,m$^{-1}$): the irrigation water has $EC_w=1.5$ dS\,m$^{-1}$. Find the leaching requirement and the gross depth for a net requirement of 450 mm.}\sol \givenline{$EC_e=2.5$ dS\,m$^{-1}$ ; $EC_w=1.5$ dS\,m$^{-1}$}\ans \par\indent We know that
\[LR=EC_w/(5EC_e-EC_w)\ti{i}\]
Substituting the values in equation \rfi{i},
\[LR=1.5/(5\times2.5-1.5)=0.136\] \ans Gross $=450/(1-LR)=450/0.864=521$ mm \ans Extra water for leaching $71$ mm \finans

\exbox{Crop example 14a: Potato}{Potato (threshold $EC_e=1.7$ dS\,m$^{-1}$): the irrigation water has $EC_w=1.5$ dS\,m$^{-1}$. Find the leaching requirement and the gross depth for a net requirement of 450 mm.}\sol \givenline{$EC_e=1.7$ dS\,m$^{-1}$ ; $EC_w=1.5$ dS\,m$^{-1}$}\ans \par\indent We know that
\[LR=EC_w/(5EC_e-EC_w)\ti{i}\]
Substituting the values in equation \rfi{i},
\[LR=1.5/(5\times1.7-1.5)=0.214\] \ans Gross $=450/(1-LR)=450/0.786=573$ mm \ans Extra water for leaching $123$ mm \finans

\exbox{Crop example 14a: Soybean}{Soybean (threshold $EC_e=5.0$ dS\,m$^{-1}$): the irrigation water has $EC_w=1.5$ dS\,m$^{-1}$. Find the leaching requirement and the gross depth for a net requirement of 450 mm.}\sol \givenline{$EC_e=5.0$ dS\,m$^{-1}$ ; $EC_w=1.5$ dS\,m$^{-1}$}\ans \par\indent We know that
\[LR=EC_w/(5EC_e-EC_w)\ti{i}\]
Substituting the values in equation \rfi{i},
\[LR=1.5/(5\times5.0-1.5)=0.064\] \ans Gross $=450/(1-LR)=450/0.936=481$ mm \ans Extra water for leaching $31$ mm \finans

\exbox{Crop example 14a: Sugar beet}{Sugar beet (threshold $EC_e=7.0$ dS\,m$^{-1}$): the irrigation water has $EC_w=1.5$ dS\,m$^{-1}$. Find the leaching requirement and the gross depth for a net requirement of 450 mm.}\sol \givenline{$EC_e=7.0$ dS\,m$^{-1}$ ; $EC_w=1.5$ dS\,m$^{-1}$}\ans \par\indent We know that
\[LR=EC_w/(5EC_e-EC_w)\ti{i}\]
Substituting the values in equation \rfi{i},
\[LR=1.5/(5\times7.0-1.5)=0.045\] \ans Gross $=450/(1-LR)=450/0.955=471$ mm \ans Extra water for leaching $21$ mm \finans

\exbox{Crop example 14a: Sunflower}{Sunflower (threshold $EC_e=4.8$ dS\,m$^{-1}$): the irrigation water has $EC_w=1.5$ dS\,m$^{-1}$. Find the leaching requirement and the gross depth for a net requirement of 450 mm.}\sol \givenline{$EC_e=4.8$ dS\,m$^{-1}$ ; $EC_w=1.5$ dS\,m$^{-1}$}\ans \par\indent We know that
\[LR=EC_w/(5EC_e-EC_w)\ti{i}\]
Substituting the values in equation \rfi{i},
\[LR=1.5/(5\times4.8-1.5)=0.067\] \ans Gross $=450/(1-LR)=450/0.933=482$ mm \ans Extra water for leaching $32$ mm \finans

\exbox{Crop example 14a: Onion}{Onion (threshold $EC_e=1.2$ dS\,m$^{-1}$): the irrigation water has $EC_w=1.5$ dS\,m$^{-1}$. Find the leaching requirement and the gross depth for a net requirement of 450 mm.}\sol \givenline{$EC_e=1.2$ dS\,m$^{-1}$ ; $EC_w=1.5$ dS\,m$^{-1}$}\ans \par\indent We know that
\[LR=EC_w/(5EC_e-EC_w)\ti{i}\]
Substituting the values in equation \rfi{i},
\[LR=1.5/(5\times1.2-1.5)=0.333\] \ans Gross $=450/(1-LR)=450/0.667=675$ mm \ans Extra water for leaching $225$ mm \finans

\exbox{Crop example 14a: Alfalfa}{Alfalfa (threshold $EC_e=2.0$ dS\,m$^{-1}$): the irrigation water has $EC_w=1.5$ dS\,m$^{-1}$. Find the leaching requirement and the gross depth for a net requirement of 450 mm.}\sol \givenline{$EC_e=2.0$ dS\,m$^{-1}$ ; $EC_w=1.5$ dS\,m$^{-1}$}\ans \par\indent We know that
\[LR=EC_w/(5EC_e-EC_w)\ti{i}\]
Substituting the values in equation \rfi{i},
\[LR=1.5/(5\times2.0-1.5)=0.176\] \ans Gross $=450/(1-LR)=450/0.824=546$ mm \ans Extra water for leaching $96$ mm \finans

}{}
\IfFileExists{extra/teny14.tex}{\subsection{Worked solutions on the ten soils (third scenario)}
The same questions are repeated for different working conditions, to show how the answer moves with them.

\exbox{Ten-soil example 14c: Coarse loam (S01)}{Coarse loam ($K_2=0.312$ m\,d$^{-1}$ from the class value, $K_1=1.2K_2$): a drain system must remove $q=10$ mm\,d$^{-1}$ with $h=0.6$ m and an equivalent depth $d=3$ m. Find the drain spacing and the length of pipe per hectare.}\sol \givenline{$K_2=0.312$ m\,d$^{-1}$ ; $K_1=1.2K_2$ ; $q=10$ mm\,d$^{-1}$ ; $h=0.6$ m ; $d=3$ m}\ans $8K_2dh=8\times0.312\times3\times0.6=4.493$; $4K_1h^2=4\times0.374\times0.36=0.539$ \ans $L^2=(4.493+0.539)/0.010=503$, $L=22$ m \ans Pipe: $10^4/L=446$ m\,ha$^{-1}$ \ans The soil is permeable: a wide spacing is possible. \finans

\exbox{Ten-soil example 14c: Coarse sand (S02)}{Coarse sand ($K_2=2.880$ m\,d$^{-1}$ from the class value, $K_1=1.2K_2$): a drain system must remove $q=10$ mm\,d$^{-1}$ with $h=0.6$ m and an equivalent depth $d=3$ m. Find the drain spacing and the length of pipe per hectare.}\sol \givenline{$K_2=2.880$ m\,d$^{-1}$ ; $K_1=1.2K_2$ ; $q=10$ mm\,d$^{-1}$ ; $h=0.6$ m ; $d=3$ m}\ans $8K_2dh=8\times2.880\times3\times0.6=41.472$; $4K_1h^2=4\times3.456\times0.36=4.977$ \ans $L^2=(41.472+4.977)/0.010=4645$, $L=68$ m \ans Pipe: $10^4/L=147$ m\,ha$^{-1}$ \ans The soil is permeable: a wide spacing is possible. \finans

\exbox{Ten-soil example 14c: Sand (S03)}{Sand ($K_2=5.040$ m\,d$^{-1}$ from the class value, $K_1=1.2K_2$): a drain system must remove $q=10$ mm\,d$^{-1}$ with $h=0.6$ m and an equivalent depth $d=3$ m. Find the drain spacing and the length of pipe per hectare.}\sol \givenline{$K_2=5.040$ m\,d$^{-1}$ ; $K_1=1.2K_2$ ; $q=10$ mm\,d$^{-1}$ ; $h=0.6$ m ; $d=3$ m}\ans $8K_2dh=8\times5.040\times3\times0.6=72.576$; $4K_1h^2=4\times6.048\times0.36=8.709$ \ans $L^2=(72.576+8.709)/0.010=8129$, $L=90$ m \ans Pipe: $10^4/L=111$ m\,ha$^{-1}$ \ans The soil is permeable: a wide spacing is possible. \finans

\exbox{Ten-soil example 14c: Loamy sand (S04)}{Loamy sand ($K_2=1.464$ m\,d$^{-1}$ from the class value, $K_1=1.2K_2$): a drain system must remove $q=10$ mm\,d$^{-1}$ with $h=0.6$ m and an equivalent depth $d=3$ m. Find the drain spacing and the length of pipe per hectare.}\sol \givenline{$K_2=1.464$ m\,d$^{-1}$ ; $K_1=1.2K_2$ ; $q=10$ mm\,d$^{-1}$ ; $h=0.6$ m ; $d=3$ m}\ans $8K_2dh=8\times1.464\times3\times0.6=21.082$; $4K_1h^2=4\times1.757\times0.36=2.530$ \ans $L^2=(21.082+2.530)/0.010=2361$, $L=49$ m \ans Pipe: $10^4/L=206$ m\,ha$^{-1}$ \ans The soil is permeable: a wide spacing is possible. \finans

\exbox{Ten-soil example 14c: Sandy loam (S05)}{Sandy loam ($K_2=0.624$ m\,d$^{-1}$ from the class value, $K_1=1.2K_2$): a drain system must remove $q=10$ mm\,d$^{-1}$ with $h=0.6$ m and an equivalent depth $d=3$ m. Find the drain spacing and the length of pipe per hectare.}\sol \givenline{$K_2=0.624$ m\,d$^{-1}$ ; $K_1=1.2K_2$ ; $q=10$ mm\,d$^{-1}$ ; $h=0.6$ m ; $d=3$ m}\ans $8K_2dh=8\times0.624\times3\times0.6=8.986$; $4K_1h^2=4\times0.749\times0.36=1.078$ \ans $L^2=(8.986+1.078)/0.010=1006$, $L=32$ m \ans Pipe: $10^4/L=315$ m\,ha$^{-1}$ \ans The soil is permeable: a wide spacing is possible. \finans

\exbox{Ten-soil example 14c: Loam (S06)}{Loam ($K_2=0.312$ m\,d$^{-1}$ from the class value, $K_1=1.2K_2$): a drain system must remove $q=10$ mm\,d$^{-1}$ with $h=0.6$ m and an equivalent depth $d=3$ m. Find the drain spacing and the length of pipe per hectare.}\sol \givenline{$K_2=0.312$ m\,d$^{-1}$ ; $K_1=1.2K_2$ ; $q=10$ mm\,d$^{-1}$ ; $h=0.6$ m ; $d=3$ m}\ans $8K_2dh=8\times0.312\times3\times0.6=4.493$; $4K_1h^2=4\times0.374\times0.36=0.539$ \ans $L^2=(4.493+0.539)/0.010=503$, $L=22$ m \ans Pipe: $10^4/L=446$ m\,ha$^{-1}$ \ans The soil is permeable: a wide spacing is possible. \finans

\exbox{Ten-soil example 14c: Silt loam (S07)}{Silt loam ($K_2=0.163$ m\,d$^{-1}$ from the class value, $K_1=1.2K_2$): a drain system must remove $q=10$ mm\,d$^{-1}$ with $h=0.6$ m and an equivalent depth $d=3$ m. Find the drain spacing and the length of pipe per hectare.}\sol \givenline{$K_2=0.163$ m\,d$^{-1}$ ; $K_1=1.2K_2$ ; $q=10$ mm\,d$^{-1}$ ; $h=0.6$ m ; $d=3$ m}\ans $8K_2dh=8\times0.163\times3\times0.6=2.350$; $4K_1h^2=4\times0.196\times0.36=0.282$ \ans $L^2=(2.350+0.282)/0.010=263$, $L=16$ m \ans Pipe: $10^4/L=616$ m\,ha$^{-1}$ \ans The soil is slow: spacing is narrow and the cost high. \finans

\exbox{Ten-soil example 14c: Sandy clay loam (S08)}{Sandy clay loam ($K_2=0.103$ m\,d$^{-1}$ from the class value, $K_1=1.2K_2$): a drain system must remove $q=10$ mm\,d$^{-1}$ with $h=0.6$ m and an equivalent depth $d=3$ m. Find the drain spacing and the length of pipe per hectare.}\sol \givenline{$K_2=0.103$ m\,d$^{-1}$ ; $K_1=1.2K_2$ ; $q=10$ mm\,d$^{-1}$ ; $h=0.6$ m ; $d=3$ m}\ans $8K_2dh=8\times0.103\times3\times0.6=1.486$; $4K_1h^2=4\times0.124\times0.36=0.178$ \ans $L^2=(1.486+0.178)/0.010=166$, $L=13$ m \ans Pipe: $10^4/L=775$ m\,ha$^{-1}$ \ans The soil is slow: spacing is narrow and the cost high. \finans

\exbox{Ten-soil example 14c: Clay loam (S09)}{Clay loam ($K_2=0.055$ m\,d$^{-1}$ from the class value, $K_1=1.2K_2$): a drain system must remove $q=10$ mm\,d$^{-1}$ with $h=0.6$ m and an equivalent depth $d=3$ m. Find the drain spacing and the length of pipe per hectare.}\sol \givenline{$K_2=0.055$ m\,d$^{-1}$ ; $K_1=1.2K_2$ ; $q=10$ mm\,d$^{-1}$ ; $h=0.6$ m ; $d=3$ m}\ans $8K_2dh=8\times0.055\times3\times0.6=0.795$; $4K_1h^2=4\times0.066\times0.36=0.095$ \ans $L^2=(0.795+0.095)/0.010=89$, $L=9$ m \ans Pipe: $10^4/L=1060$ m\,ha$^{-1}$ \ans The soil is slow: spacing is narrow and the cost high. \finans

\exbox{Ten-soil example 14c: Silty clay loam (S10)}{Silty clay loam ($K_2=0.036$ m\,d$^{-1}$ from the class value, $K_1=1.2K_2$): a drain system must remove $q=10$ mm\,d$^{-1}$ with $h=0.6$ m and an equivalent depth $d=3$ m. Find the drain spacing and the length of pipe per hectare.}\sol \givenline{$K_2=0.036$ m\,d$^{-1}$ ; $K_1=1.2K_2$ ; $q=10$ mm\,d$^{-1}$ ; $h=0.6$ m ; $d=3$ m}\ans $8K_2dh=8\times0.036\times3\times0.6=0.518$; $4K_1h^2=4\times0.043\times0.36=0.062$ \ans $L^2=(0.518+0.062)/0.010=58$, $L=8$ m \ans Pipe: $10^4/L=1312$ m\,ha$^{-1}$ \ans The soil is slow: spacing is narrow and the cost high. \finans

}{}
\IfFileExists{extra/tenby14.tex}{\section{Worked solutions on ten crops (hot climate)}
The ten crops of the library differ in their demand, their roots and their sensitivity. The same type of question is answered for each of them at $ET_0=7.5$ mm\,d$^{-1}$; the data are typical values of the crops (Doorenbos--Kassam and FAO 56).

\exbox{Crop example 14c: Maize}{Maize (threshold $EC_e=1.7$ dS\,m$^{-1}$): the irrigation water has $EC_w=1.5$ dS\,m$^{-1}$. Find the leaching requirement and the gross depth for a net requirement of 450 mm.}\sol $LR=EC_w/(5EC_e-EC_w)=1.5/(5\times1.7-1.5)=0.214$ \ans Gross $=450/(1-LR)=450/0.786=573$ mm \ans Extra water for leaching $123$ mm \ans

\exbox{Crop example 14c: Wheat}{Wheat (threshold $EC_e=6.0$ dS\,m$^{-1}$): the irrigation water has $EC_w=1.5$ dS\,m$^{-1}$. Find the leaching requirement and the gross depth for a net requirement of 450 mm.}\sol $LR=EC_w/(5EC_e-EC_w)=1.5/(5\times6.0-1.5)=0.053$ \ans Gross $=450/(1-LR)=450/0.947=475$ mm \ans Extra water for leaching $25$ mm \ans

\exbox{Crop example 14c: Cotton}{Cotton (threshold $EC_e=7.7$ dS\,m$^{-1}$): the irrigation water has $EC_w=1.5$ dS\,m$^{-1}$. Find the leaching requirement and the gross depth for a net requirement of 450 mm.}\sol $LR=EC_w/(5EC_e-EC_w)=1.5/(5\times7.7-1.5)=0.041$ \ans Gross $=450/(1-LR)=450/0.959=469$ mm \ans Extra water for leaching $19$ mm \ans

\exbox{Crop example 14c: Tomato}{Tomato (threshold $EC_e=2.5$ dS\,m$^{-1}$): the irrigation water has $EC_w=1.5$ dS\,m$^{-1}$. Find the leaching requirement and the gross depth for a net requirement of 450 mm.}\sol $LR=EC_w/(5EC_e-EC_w)=1.5/(5\times2.5-1.5)=0.136$ \ans Gross $=450/(1-LR)=450/0.864=521$ mm \ans Extra water for leaching $71$ mm \ans

\exbox{Crop example 14c: Potato}{Potato (threshold $EC_e=1.7$ dS\,m$^{-1}$): the irrigation water has $EC_w=1.5$ dS\,m$^{-1}$. Find the leaching requirement and the gross depth for a net requirement of 450 mm.}\sol $LR=EC_w/(5EC_e-EC_w)=1.5/(5\times1.7-1.5)=0.214$ \ans Gross $=450/(1-LR)=450/0.786=573$ mm \ans Extra water for leaching $123$ mm \ans

\exbox{Crop example 14c: Soybean}{Soybean (threshold $EC_e=5.0$ dS\,m$^{-1}$): the irrigation water has $EC_w=1.5$ dS\,m$^{-1}$. Find the leaching requirement and the gross depth for a net requirement of 450 mm.}\sol $LR=EC_w/(5EC_e-EC_w)=1.5/(5\times5.0-1.5)=0.064$ \ans Gross $=450/(1-LR)=450/0.936=481$ mm \ans Extra water for leaching $31$ mm \ans

\exbox{Crop example 14c: Sugar beet}{Sugar beet (threshold $EC_e=7.0$ dS\,m$^{-1}$): the irrigation water has $EC_w=1.5$ dS\,m$^{-1}$. Find the leaching requirement and the gross depth for a net requirement of 450 mm.}\sol $LR=EC_w/(5EC_e-EC_w)=1.5/(5\times7.0-1.5)=0.045$ \ans Gross $=450/(1-LR)=450/0.955=471$ mm \ans Extra water for leaching $21$ mm \ans

\exbox{Crop example 14c: Sunflower}{Sunflower (threshold $EC_e=4.8$ dS\,m$^{-1}$): the irrigation water has $EC_w=1.5$ dS\,m$^{-1}$. Find the leaching requirement and the gross depth for a net requirement of 450 mm.}\sol $LR=EC_w/(5EC_e-EC_w)=1.5/(5\times4.8-1.5)=0.067$ \ans Gross $=450/(1-LR)=450/0.933=482$ mm \ans Extra water for leaching $32$ mm \ans

\exbox{Crop example 14c: Onion}{Onion (threshold $EC_e=1.2$ dS\,m$^{-1}$): the irrigation water has $EC_w=1.5$ dS\,m$^{-1}$. Find the leaching requirement and the gross depth for a net requirement of 450 mm.}\sol $LR=EC_w/(5EC_e-EC_w)=1.5/(5\times1.2-1.5)=0.333$ \ans Gross $=450/(1-LR)=450/0.667=675$ mm \ans Extra water for leaching $225$ mm \ans

\exbox{Crop example 14c: Alfalfa}{Alfalfa (threshold $EC_e=2.0$ dS\,m$^{-1}$): the irrigation water has $EC_w=1.5$ dS\,m$^{-1}$. Find the leaching requirement and the gross depth for a net requirement of 450 mm.}\sol $LR=EC_w/(5EC_e-EC_w)=1.5/(5\times2.0-1.5)=0.176$ \ans Gross $=450/(1-LR)=450/0.824=546$ mm \ans Extra water for leaching $96$ mm \ans

}{}
\IfFileExists{extra/tenx14.tex}{\subsection{Worked solutions on the ten soils (second scenario)}
The same questions are repeated for different working conditions, to show how the answer moves with them.

\exbox{Ten-soil example 14b: Coarse loam (S01)}{Coarse loam ($K_2=0.312$ m\,d$^{-1}$ from the class value, $K_1=1.2K_2$): a drain system must remove $q=4$ mm\,d$^{-1}$ with $h=0.4$ m and an equivalent depth $d=3$ m. Find the drain spacing and the length of pipe per hectare.}\sol \givenline{$K_2=0.312$ m\,d$^{-1}$ ; $K_1=1.2K_2$ ; $q=4$ mm\,d$^{-1}$ ; $h=0.4$ m ; $d=3$ m}\ans \[8K_2dh=8\times0.312\times3\times0.4=2.995\]
\[4K_1h^2=4\times0.374\times0.16=0.240\] \ans \[L^2=(2.995+0.240)/0.004=809\]
$L=28$ m \ans Pipe: $10^4/L=352$ m\,ha$^{-1}$ \ans The soil is permeable: a wide spacing is possible. \finans

\exbox{Ten-soil example 14b: Coarse sand (S02)}{Coarse sand ($K_2=2.880$ m\,d$^{-1}$ from the class value, $K_1=1.2K_2$): a drain system must remove $q=4$ mm\,d$^{-1}$ with $h=0.4$ m and an equivalent depth $d=3$ m. Find the drain spacing and the length of pipe per hectare.}\sol \givenline{$K_2=2.880$ m\,d$^{-1}$ ; $K_1=1.2K_2$ ; $q=4$ mm\,d$^{-1}$ ; $h=0.4$ m ; $d=3$ m}\ans \[8K_2dh=8\times2.880\times3\times0.4=27.648\]
\[4K_1h^2=4\times3.456\times0.16=2.212\] \ans \[L^2=(27.648+2.212)/0.004=7465\]
$L=86$ m \ans Pipe: $10^4/L=116$ m\,ha$^{-1}$ \ans The soil is permeable: a wide spacing is possible. \finans

\exbox{Ten-soil example 14b: Sand (S03)}{Sand ($K_2=5.040$ m\,d$^{-1}$ from the class value, $K_1=1.2K_2$): a drain system must remove $q=4$ mm\,d$^{-1}$ with $h=0.4$ m and an equivalent depth $d=3$ m. Find the drain spacing and the length of pipe per hectare.}\sol \givenline{$K_2=5.040$ m\,d$^{-1}$ ; $K_1=1.2K_2$ ; $q=4$ mm\,d$^{-1}$ ; $h=0.4$ m ; $d=3$ m}\ans \[8K_2dh=8\times5.040\times3\times0.4=48.384\]
\[4K_1h^2=4\times6.048\times0.16=3.871\] \ans \[L^2=(48.384+3.871)/0.004=13064\]
$L=114$ m \ans Pipe: $10^4/L=87$ m\,ha$^{-1}$ \ans The soil is permeable: a wide spacing is possible. \finans

\exbox{Ten-soil example 14b: Loamy sand (S04)}{Loamy sand ($K_2=1.464$ m\,d$^{-1}$ from the class value, $K_1=1.2K_2$): a drain system must remove $q=4$ mm\,d$^{-1}$ with $h=0.4$ m and an equivalent depth $d=3$ m. Find the drain spacing and the length of pipe per hectare.}\sol \givenline{$K_2=1.464$ m\,d$^{-1}$ ; $K_1=1.2K_2$ ; $q=4$ mm\,d$^{-1}$ ; $h=0.4$ m ; $d=3$ m}\ans \[8K_2dh=8\times1.464\times3\times0.4=14.054\]
\[4K_1h^2=4\times1.757\times0.16=1.124\] \ans \[L^2=(14.054+1.124)/0.004=3795\]
$L=62$ m \ans Pipe: $10^4/L=162$ m\,ha$^{-1}$ \ans The soil is permeable: a wide spacing is possible. \finans

\exbox{Ten-soil example 14b: Sandy loam (S05)}{Sandy loam ($K_2=0.624$ m\,d$^{-1}$ from the class value, $K_1=1.2K_2$): a drain system must remove $q=4$ mm\,d$^{-1}$ with $h=0.4$ m and an equivalent depth $d=3$ m. Find the drain spacing and the length of pipe per hectare.}\sol \givenline{$K_2=0.624$ m\,d$^{-1}$ ; $K_1=1.2K_2$ ; $q=4$ mm\,d$^{-1}$ ; $h=0.4$ m ; $d=3$ m}\ans \[8K_2dh=8\times0.624\times3\times0.4=5.990\]
\[4K_1h^2=4\times0.749\times0.16=0.479\] \ans \[L^2=(5.990+0.479)/0.004=1617\]
$L=40$ m \ans Pipe: $10^4/L=249$ m\,ha$^{-1}$ \ans The soil is permeable: a wide spacing is possible. \finans

\exbox{Ten-soil example 14b: Loam (S06)}{Loam ($K_2=0.312$ m\,d$^{-1}$ from the class value, $K_1=1.2K_2$): a drain system must remove $q=4$ mm\,d$^{-1}$ with $h=0.4$ m and an equivalent depth $d=3$ m. Find the drain spacing and the length of pipe per hectare.}\sol \givenline{$K_2=0.312$ m\,d$^{-1}$ ; $K_1=1.2K_2$ ; $q=4$ mm\,d$^{-1}$ ; $h=0.4$ m ; $d=3$ m}\ans \[8K_2dh=8\times0.312\times3\times0.4=2.995\]
\[4K_1h^2=4\times0.374\times0.16=0.240\] \ans \[L^2=(2.995+0.240)/0.004=809\]
$L=28$ m \ans Pipe: $10^4/L=352$ m\,ha$^{-1}$ \ans The soil is permeable: a wide spacing is possible. \finans

\exbox{Ten-soil example 14b: Silt loam (S07)}{Silt loam ($K_2=0.163$ m\,d$^{-1}$ from the class value, $K_1=1.2K_2$): a drain system must remove $q=4$ mm\,d$^{-1}$ with $h=0.4$ m and an equivalent depth $d=3$ m. Find the drain spacing and the length of pipe per hectare.}\sol \givenline{$K_2=0.163$ m\,d$^{-1}$ ; $K_1=1.2K_2$ ; $q=4$ mm\,d$^{-1}$ ; $h=0.4$ m ; $d=3$ m}\ans \[8K_2dh=8\times0.163\times3\times0.4=1.567\]
\[4K_1h^2=4\times0.196\times0.16=0.125\] \ans \[L^2=(1.567+0.125)/0.004=423\]
$L=21$ m \ans Pipe: $10^4/L=486$ m\,ha$^{-1}$ \ans The soil is permeable: a wide spacing is possible. \finans

\exbox{Ten-soil example 14b: Sandy clay loam (S08)}{Sandy clay loam ($K_2=0.103$ m\,d$^{-1}$ from the class value, $K_1=1.2K_2$): a drain system must remove $q=4$ mm\,d$^{-1}$ with $h=0.4$ m and an equivalent depth $d=3$ m. Find the drain spacing and the length of pipe per hectare.}\sol \givenline{$K_2=0.103$ m\,d$^{-1}$ ; $K_1=1.2K_2$ ; $q=4$ mm\,d$^{-1}$ ; $h=0.4$ m ; $d=3$ m}\ans \[8K_2dh=8\times0.103\times3\times0.4=0.991\]
\[4K_1h^2=4\times0.124\times0.16=0.079\] \ans \[L^2=(0.991+0.079)/0.004=267\]
$L=16$ m \ans Pipe: $10^4/L=611$ m\,ha$^{-1}$ \ans The soil is slow: spacing is narrow and the cost high. \finans

\exbox{Ten-soil example 14b: Clay loam (S09)}{Clay loam ($K_2=0.055$ m\,d$^{-1}$ from the class value, $K_1=1.2K_2$): a drain system must remove $q=4$ mm\,d$^{-1}$ with $h=0.4$ m and an equivalent depth $d=3$ m. Find the drain spacing and the length of pipe per hectare.}\sol \givenline{$K_2=0.055$ m\,d$^{-1}$ ; $K_1=1.2K_2$ ; $q=4$ mm\,d$^{-1}$ ; $h=0.4$ m ; $d=3$ m}\ans \[8K_2dh=8\times0.055\times3\times0.4=0.530\]
\[4K_1h^2=4\times0.066\times0.16=0.042\] \ans \[L^2=(0.530+0.042)/0.004=143\]
$L=12$ m \ans Pipe: $10^4/L=836$ m\,ha$^{-1}$ \ans The soil is slow: spacing is narrow and the cost high. \finans

\exbox{Ten-soil example 14b: Silty clay loam (S10)}{Silty clay loam ($K_2=0.036$ m\,d$^{-1}$ from the class value, $K_1=1.2K_2$): a drain system must remove $q=4$ mm\,d$^{-1}$ with $h=0.4$ m and an equivalent depth $d=3$ m. Find the drain spacing and the length of pipe per hectare.}\sol \givenline{$K_2=0.036$ m\,d$^{-1}$ ; $K_1=1.2K_2$ ; $q=4$ mm\,d$^{-1}$ ; $h=0.4$ m ; $d=3$ m}\ans \[8K_2dh=8\times0.036\times3\times0.4=0.346\]
\[4K_1h^2=4\times0.043\times0.16=0.028\] \ans \[L^2=(0.346+0.028)/0.004=93\]
$L=10$ m \ans Pipe: $10^4/L=1035$ m\,ha$^{-1}$ \ans The soil is slow: spacing is narrow and the cost high. \finans

}{}
\IfFileExists{extra/tenbx14.tex}{\subsection{Worked solutions on ten crops (cool climate)}
The ten crops of the library differ in their demand, their roots and their sensitivity. The same type of question is answered for each of them at $ET_0=4.0$ mm\,d$^{-1}$; the data are typical values of the crops (Doorenbos--Kassam and FAO 56).

\exbox{Crop example 14b: Maize}{Maize (threshold $EC_e=1.7$ dS\,m$^{-1}$): the irrigation water has $EC_w=1.5$ dS\,m$^{-1}$. Find the leaching requirement and the gross depth for a net requirement of 450 mm.}\sol \givenline{$EC_e=1.7$ dS\,m$^{-1}$ ; $EC_w=1.5$ dS\,m$^{-1}$}\ans $LR=EC_w/(5EC_e-EC_w)=1.5/(5\times1.7-1.5)=0.214$ \ans Gross $=450/(1-LR)=450/0.786=573$ mm \ans Extra water for leaching $123$ mm \finans

\exbox{Crop example 14b: Wheat}{Wheat (threshold $EC_e=6.0$ dS\,m$^{-1}$): the irrigation water has $EC_w=1.5$ dS\,m$^{-1}$. Find the leaching requirement and the gross depth for a net requirement of 450 mm.}\sol \givenline{$EC_e=6.0$ dS\,m$^{-1}$ ; $EC_w=1.5$ dS\,m$^{-1}$}\ans $LR=EC_w/(5EC_e-EC_w)=1.5/(5\times6.0-1.5)=0.053$ \ans Gross $=450/(1-LR)=450/0.947=475$ mm \ans Extra water for leaching $25$ mm \finans

\exbox{Crop example 14b: Cotton}{Cotton (threshold $EC_e=7.7$ dS\,m$^{-1}$): the irrigation water has $EC_w=1.5$ dS\,m$^{-1}$. Find the leaching requirement and the gross depth for a net requirement of 450 mm.}\sol \givenline{$EC_e=7.7$ dS\,m$^{-1}$ ; $EC_w=1.5$ dS\,m$^{-1}$}\ans $LR=EC_w/(5EC_e-EC_w)=1.5/(5\times7.7-1.5)=0.041$ \ans Gross $=450/(1-LR)=450/0.959=469$ mm \ans Extra water for leaching $19$ mm \finans

\exbox{Crop example 14b: Tomato}{Tomato (threshold $EC_e=2.5$ dS\,m$^{-1}$): the irrigation water has $EC_w=1.5$ dS\,m$^{-1}$. Find the leaching requirement and the gross depth for a net requirement of 450 mm.}\sol \givenline{$EC_e=2.5$ dS\,m$^{-1}$ ; $EC_w=1.5$ dS\,m$^{-1}$}\ans $LR=EC_w/(5EC_e-EC_w)=1.5/(5\times2.5-1.5)=0.136$ \ans Gross $=450/(1-LR)=450/0.864=521$ mm \ans Extra water for leaching $71$ mm \finans

\exbox{Crop example 14b: Potato}{Potato (threshold $EC_e=1.7$ dS\,m$^{-1}$): the irrigation water has $EC_w=1.5$ dS\,m$^{-1}$. Find the leaching requirement and the gross depth for a net requirement of 450 mm.}\sol \givenline{$EC_e=1.7$ dS\,m$^{-1}$ ; $EC_w=1.5$ dS\,m$^{-1}$}\ans $LR=EC_w/(5EC_e-EC_w)=1.5/(5\times1.7-1.5)=0.214$ \ans Gross $=450/(1-LR)=450/0.786=573$ mm \ans Extra water for leaching $123$ mm \finans

\exbox{Crop example 14b: Soybean}{Soybean (threshold $EC_e=5.0$ dS\,m$^{-1}$): the irrigation water has $EC_w=1.5$ dS\,m$^{-1}$. Find the leaching requirement and the gross depth for a net requirement of 450 mm.}\sol \givenline{$EC_e=5.0$ dS\,m$^{-1}$ ; $EC_w=1.5$ dS\,m$^{-1}$}\ans $LR=EC_w/(5EC_e-EC_w)=1.5/(5\times5.0-1.5)=0.064$ \ans Gross $=450/(1-LR)=450/0.936=481$ mm \ans Extra water for leaching $31$ mm \finans

\exbox{Crop example 14b: Sugar beet}{Sugar beet (threshold $EC_e=7.0$ dS\,m$^{-1}$): the irrigation water has $EC_w=1.5$ dS\,m$^{-1}$. Find the leaching requirement and the gross depth for a net requirement of 450 mm.}\sol \givenline{$EC_e=7.0$ dS\,m$^{-1}$ ; $EC_w=1.5$ dS\,m$^{-1}$}\ans $LR=EC_w/(5EC_e-EC_w)=1.5/(5\times7.0-1.5)=0.045$ \ans Gross $=450/(1-LR)=450/0.955=471$ mm \ans Extra water for leaching $21$ mm \finans

\exbox{Crop example 14b: Sunflower}{Sunflower (threshold $EC_e=4.8$ dS\,m$^{-1}$): the irrigation water has $EC_w=1.5$ dS\,m$^{-1}$. Find the leaching requirement and the gross depth for a net requirement of 450 mm.}\sol \givenline{$EC_e=4.8$ dS\,m$^{-1}$ ; $EC_w=1.5$ dS\,m$^{-1}$}\ans $LR=EC_w/(5EC_e-EC_w)=1.5/(5\times4.8-1.5)=0.067$ \ans Gross $=450/(1-LR)=450/0.933=482$ mm \ans Extra water for leaching $32$ mm \finans

\exbox{Crop example 14b: Onion}{Onion (threshold $EC_e=1.2$ dS\,m$^{-1}$): the irrigation water has $EC_w=1.5$ dS\,m$^{-1}$. Find the leaching requirement and the gross depth for a net requirement of 450 mm.}\sol \givenline{$EC_e=1.2$ dS\,m$^{-1}$ ; $EC_w=1.5$ dS\,m$^{-1}$}\ans $LR=EC_w/(5EC_e-EC_w)=1.5/(5\times1.2-1.5)=0.333$ \ans Gross $=450/(1-LR)=450/0.667=675$ mm \ans Extra water for leaching $225$ mm \finans

\exbox{Crop example 14b: Alfalfa}{Alfalfa (threshold $EC_e=2.0$ dS\,m$^{-1}$): the irrigation water has $EC_w=1.5$ dS\,m$^{-1}$. Find the leaching requirement and the gross depth for a net requirement of 450 mm.}\sol \givenline{$EC_e=2.0$ dS\,m$^{-1}$ ; $EC_w=1.5$ dS\,m$^{-1}$}\ans $LR=EC_w/(5EC_e-EC_w)=1.5/(5\times2.0-1.5)=0.176$ \ans Gross $=450/(1-LR)=450/0.824=546$ mm \ans Extra water for leaching $96$ mm \finans

}{}
\IfFileExists{extra/ex14.tex}{\backhead{EXERCISES}{khred}
\begin{exlist}
\item $LR$ for $EC_w=2$, $EC_e=6$.\exans{0.071}
\item $q$ for $L=40$, $K_2=0.5$, $d_e=2$, $h=0.4$, $K_1=0.8$.\exans{0.0023 m/d}
\item Spacing for $q=5$ mm/d if $8K_2d_eh+4K_1h^2=20$.\exans{63 m}
\item Pipe per hectare for $L=25$ m.\exans{400 m}
\item Leaching depth for 800 mm applied, LR 0.15.\exans{120 mm}
\item Why is $d_e$ less than $d$?\exans{flow convergence}
\item Name two kinds of subsurface drains.\exans{pipe; mole}
\item What is $r_0$?\exans{drain radius}
\item How does a barrier affect $L$?\exans{shallower barrier, narrower spacing}
\item Why leach?\exans{remove salt}
\end{exlist}
}{}
\IfFileExists{extra/dyk14.tex}{\dykband{DO YOU KNOW ?}
\begin{exlist}
\item Name the types of drainage.
\item What does the Hooghoudt equation give?
\item Why does $d_e$ depend on $L$?
\item What is the leaching requirement?
\item How is the spacing solved?
\item What limits drain depth?
\item What is a mole drain?
\item Why do salinity and drainage go together?
\end{exlist}
}{}
\IfFileExists{extra/obj14.tex}{\dykband{OBJECTIVE TYPE QUESTIONS}
\begin{exlist}
\item Hooghoudt gives\opt{spacing}{rain}{ET}{price}
\item $LR=$\opt{$EC_w/(5EC_e-EC_w)$}{$EC_e/EC_w$}{$EC_w$}{none}
\item $d_e$ is\opt{equivalent depth}{drain size}{rain}{none}
\item Barrier shallow\opt{narrower spacing}{wider}{none}{salt}
\item Mole drain is\opt{unlined channel}{pipe}{pump}{none}
\item $K_1$ is\opt{above drains}{below}{none}{salt}
\item Leaching removes\opt{salt}{rain}{crop}{none}
\item Spacing is by\opt{iteration}{guess}{none}{rain}
\item $r_0$ is\opt{drain radius}{rain}{salt}{none}
\item $h$ is\opt{head above drain}{depth}{time}{none}
\end{exlist}
\objkey{1.~(a)\quad 2.~(a)\quad 3.~(a)\quad 4.~(a)\quad 5.~(a)\quad 6.~(a)\quad 7.~(a)\quad 8.~(a)\quad 9.~(a)\quad 10.~(a)}
}{}
\swaflvideo{14}
